% ============================================================
%  Manual Técnico — K-GeoMIP
%  Proyecto: K-GeoMIP — Extensión de GeoMIP a k-Particiones
%  Curso:    Análisis y Diseño de Algoritmos — 2026-1
%  Autor:    Dylber Cabrera
%  Fecha:    2026-06-15
% ============================================================
\documentclass[12pt,letterpaper]{article}

% ── Codificación y lenguaje ─────────────────────────────────
\usepackage[utf8]{inputenc}
\usepackage[T1]{fontenc}
\usepackage[spanish,es-tabla]{babel}
\usepackage{csquotes}

% ── Diseño de página ────────────────────────────────────────
\usepackage[left=2.5cm,right=2.5cm,top=2.5cm,bottom=2.5cm,
            headheight=14pt]{geometry}
\usepackage{fancyhdr}
\usepackage{setspace}
\usepackage{parskip}

% ── Tipografía ──────────────────────────────────────────────
\usepackage{lmodern}
\usepackage{microtype}

% ── Matemáticas ─────────────────────────────────────────────
\usepackage{amsmath}
\usepackage{amssymb}
\usepackage{amsthm}
\usepackage{mathtools}
\usepackage{bm}

% ── Colores y tablas ────────────────────────────────────────
\usepackage[dvipsnames,svgnames,x11names]{xcolor}
\usepackage{colortbl}  % \cellcolor, \rowcolor inside tabular

\definecolor{KGeoBlue}{RGB}{0,102,204}
\definecolor{KGeoGreen}{RGB}{34,139,34}
\definecolor{KGeoOrange}{RGB}{210,105,30}
\definecolor{KGeoGray}{RGB}{128,128,128}
\definecolor{codebg}{RGB}{248,248,248}
\definecolor{codeframe}{RGB}{180,180,180}
\definecolor{pseudobg}{RGB}{240,248,255}

% ── Gráficos y diagramas ────────────────────────────────────
\usepackage{float}      % especificador [H] para posición exacta de floats
\usepackage{graphicx}
\usepackage{tikz}
\usetikzlibrary{shapes,arrows.meta,positioning,fit,calc,
                backgrounds,decorations.pathreplacing}
\usepackage{pgfplots}
\pgfplotsset{compat=1.18}

% ── Tablas ──────────────────────────────────────────────────
\usepackage{booktabs}
\usepackage{longtable}
\usepackage{tabularx}
\usepackage{multirow}
\usepackage{array}
\usepackage{makecell}

% ── Código fuente ───────────────────────────────────────────
\usepackage{listings}
\lstdefinestyle{python}{
  language=Python,
  backgroundcolor=\color{codebg},
  frame=single,
  rulecolor=\color{codeframe},
  basicstyle=\ttfamily\small,
  keywordstyle=\color{KGeoBlue}\bfseries,
  stringstyle=\color{KGeoGreen},
  commentstyle=\color{KGeoGray}\itshape,
  numberstyle=\tiny\color{KGeoGray},
  numbers=left,
  numbersep=6pt,
  breaklines=true,
  captionpos=b,
  showstringspaces=false,
  tabsize=4,
  literate=
    {->}{{$\rightarrow$}}{2}
    {<-}{{$\leftarrow$}}{2},
  morekeywords={yield,from,self,True,False,None,def,class,
    import,return,for,if,in,not,and,or,else,elif,
    while,continue,break,pass},
}
\lstdefinestyle{pseudocode}{
  backgroundcolor=\color{pseudobg},
  frame=single,
  rulecolor=\color{KGeoBlue!40},
  basicstyle=\ttfamily\small,
  keywordstyle=\color{KGeoBlue}\bfseries,
  commentstyle=\color{KGeoGray}\itshape,
  breaklines=true,
  tabsize=2,
  morekeywords={FUNCION,ENTRADA,SALIDA,PARA,CADA,EN,SI,SINO,
    RETORNAR,EMITIR,VALIDAR,DESDE,HASTA,INTENTAR,EXCEPCION,
    IMPRIMIR,MIENTRAS,HACER},
}

% ── Algoritmos ──────────────────────────────────────────────
\usepackage[ruled,vlined,linesnumbered,spanish]{algorithm2e}
\SetAlgoCaptionSeparator{.}

% ── Hiperenlaces ────────────────────────────────────────────
\usepackage{hyperref}
\hypersetup{
  colorlinks=true,
  linkcolor=black,
  citecolor=KGeoGreen,
  urlcolor=KGeoOrange,
  pdfauthor={Dylber Cabrera},
  pdftitle={Manual Tecnico -- K-GeoMIP},
  pdfsubject={Analisis y Diseno de Algoritmos 2026-1},
}

% ── Entornos de teorema ─────────────────────────────────────
\theoremstyle{definition}
\newtheorem{definition}{Definición}[section]
\newtheorem{proposition}{Proposición}[section]
\newtheorem{lemma}{Lema}[section]
\newtheorem{corollary}{Corolario}[section]
\theoremstyle{remark}
\newtheorem{remark}{Observación}[section]
\newtheorem{example}{Ejemplo}[section]

% ── Cajas con mdframed ──────────────────────────────────────
\usepackage{mdframed}
\newmdenv[
  linecolor=KGeoBlue,
  backgroundcolor=KGeoBlue!5,
  linewidth=1.5pt,
  innerleftmargin=10pt,innerrightmargin=10pt,
  innertopmargin=8pt,innerbottommargin=8pt,
]{blueBox}
\newmdenv[
  linecolor=KGeoOrange,
  backgroundcolor=KGeoOrange!5,
  linewidth=1.5pt,
  innerleftmargin=10pt,innerrightmargin=10pt,
  innertopmargin=8pt,innerbottommargin=8pt,
]{orangeBox}

% ── Encabezados ─────────────────────────────────────────────
\pagestyle{fancy}
\fancyhf{}
\fancyhead[L]{\small Manual Técnico --- K-GeoMIP}
\fancyhead[R]{\small Análisis y Diseño de Algoritmos 2026-1}
\fancyfoot[C]{\thepage}
\renewcommand{\headrulewidth}{0.4pt}

% ── Comandos matemáticos propios ────────────────────────────
\DeclareMathOperator*{\argmin}{arg\,min}
\newcommand{\V}{\mathcal{V}}
\newcommand{\Part}{\mathcal{P}}
\newcommand{\EMD}{\operatorname{EMD}}
\newcommand{\Stir}[2]{S(#1,#2)}
\newcommand{\BigO}[1]{\mathcal{O}(#1)}
\newcommand{\BigTheta}[1]{\Theta(#1)}
\newcommand{\kMIP}{k\text{-MIP}}
\newcommand{\R}{\mathbb{R}}
\newcommand{\N}{\mathbb{N}}
\newcommand{\tensor}{\otimes}

% ── Desactivar "shorthand" de babel para evitar conflictos ──
\AtBeginDocument{\shorthandoff{"}}

% ════════════════════════════════════════════════════════════
%  PORTADA
% ════════════════════════════════════════════════════════════
\begin{document}

\begin{titlepage}
  \centering
  \vspace*{1.5cm}
  {\Large\bfseries Universidad de Caldas\par}
  {\large Facultad de Inteligencia Artificial e Ingenierías\par}
  \vspace{0.5cm}
  \rule{\linewidth}{1pt}
  \vspace{1cm}

  {\Huge\bfseries\color{KGeoBlue} Manual Técnico\par}
  \vspace{0.4cm}
  {\LARGE\bfseries K-GeoMIP\par}
  \vspace{0.3cm}
  {\large Extensión de GeoMIP a k-Particiones\par}
  \vspace{1cm}
  \rule{\linewidth}{0.5pt}
  \vspace{1cm}

  {\large
    \begin{tabular}{rl}
      \textbf{Proyecto:} & Algoritmo Geométrico K-GeoMIP \\[4pt]
      \textbf{Curso:}    & Análisis y Diseño de Algoritmos \\[4pt]
      \textbf{Período:}  & 2026-1 \\[4pt]
      \textbf{Autor:}    & Dylber Cabrera - Keyla Cartagena \\[4pt]
      \textbf{Fecha:}    & 15 de junio de 2026 \\[4pt]
      \textbf{Versión:}  & 1.0 \\
    \end{tabular}
  }
  \vfill
\end{titlepage}

\tableofcontents
\newpage

% ════════════════════════════════════════════════════════════
%  SECCIÓN 1 — RESUMEN EJECUTIVO
% ════════════════════════════════════════════════════════════
\section{Resumen Ejecutivo}

\subsection{Descripción del problema}
Este manual técnico documenta el diseño, implementación y evaluación experimental de
\textbf{K-GeoMIP}, una extensión del framework geométrico GeoMIP orientada a resolver el
problema de la \emph{k-Partición de Mínima Información} (k-MIP) en el contexto de la
Teoría de la Información Integrada (IIT). Mientras los trabajos anteriores (GeoMIP,
QNodes) se limitaban al caso de bi-particiones ($k = 2$), K-GeoMIP extiende el marco a
$k \in \{2, 3, 4, 5\}$, permitiendo analizar divisiones del sistema en múltiples partes
independientes.

\subsection{Enfoque algorítmico}
El núcleo de K-GeoMIP es un \textbf{algoritmo voraz con garantía LPT} (Longest Processing
Time first) que reduce el espacio de búsqueda de $\Stir{n}{k}^2$ evaluaciones a
$C = 25$ candidatos por llamada. La fase computacionalmente dominante, la construcción de
la tabla de costos mediante recorrido BFS sobre el hipercubo $n$-dimensional, tiene
complejidad $\BigTheta{n \cdot 2^n}$, idéntica a la del GeoMIP original para $k = 2$.

\subsection{Principales resultados}
\begin{itemize}
  \item \textbf{Sin regresión temporal:} KGeometricSIA$(k \geq 3)$ es comparable o más
    rápido que GeometricSIA$(k=2)$ en todos los datasets probados (ratio máx.\ $1{,}35\times$).
  \item \textbf{Speedup masivo en búsqueda:} Para $n=10$, $k=3$: speedup de búsqueda
    $\approx 3\,110\times$ respecto a la enumeración exhaustiva de Stirling.
  \item \textbf{Validación formal:} Suite de 13 tests automatizados confirman equivalencia
    exacta con GeoMIP para $k=2$
    ($|\delta_2^{\mathrm{KGeo}} - \delta_2^{\mathrm{Geo}}| < 10^{-9}$).
\end{itemize}

\subsection{Limitaciones}
La heurística greedy no garantiza optimalidad global para $k \geq 3$. Para sistemas con
$n > 12$ y $k > 2$, la búsqueda exhaustiva es computacionalmente inviable y la heurística
puede omitir la partición óptima. Las mejoras propuestas se documentan en la
Sección~\ref{sec:limitaciones}.

% ════════════════════════════════════════════════════════════
%  SECCIÓN 2 — FUNDAMENTOS TEÓRICOS
% ════════════════════════════════════════════════════════════
\section{Fundamentos Teóricos de k-Particiones}

\subsection{Definición formal de k-particiones}

\begin{definition}[k-Partición]
Sea $\V = \{v_1, v_2, \ldots, v_n\}$ un conjunto de $n$ variables binarias. Una
\textbf{k-partición} de $\V$ es una colección $\Part_k = \{S_1, S_2, \ldots, S_k\}$ que
satisface:
\begin{align}
  &\text{(Completitud)}    & S_1 \cup \cdots \cup S_k &= \V \label{eq:completitud}\\
  &\text{(Disjunción)}     & S_i \cap S_j &= \emptyset \quad \forall\, i \neq j
                                                        \label{eq:disjuncion}\\
  &\text{(No trivialidad)} & S_i &\neq \emptyset \quad \forall\, i \label{eq:notrivial}
\end{align}
\end{definition}

En el contexto de la IIT, cada parte $S_i$ tiene dos componentes:
$A_i$ (\textbf{alcance futuro}) y $M_i$ (\textbf{mecanismo presente}).
La k-partición del subsistema se representa como:
\begin{equation}
  \Part_k = \bigl\{ (A_1, M_1),\, (A_2, M_2),\, \ldots,\, (A_k, M_k) \bigr\}
  \label{eq:kpart-subsistema}
\end{equation}

\subsection{Formulación del problema de optimización}

\begin{definition}[k-MIP]
Dado un sistema $\V$ y Matriz de Probabilidad de Transición
$\mathrm{TPM} = P(\V_{t+1} \mid \V_t)$, la
\textbf{k-Partición de Mínima Información} (k-MIP) se define como:
\begin{equation}
  \kMIP = \argmin_{\Part \in \Pi_k(\V)} \delta_k(\Part)
  \label{eq:kmip}
\end{equation}
donde $\Pi_k(\V)$ es el conjunto de todas las k-particiones de $\V$.
\end{definition}

\subsection{Función de pérdida \texorpdfstring{$\delta_k$}{delta\_k}}

La pérdida de información para una k-partición $\Part_k$ desde el estado $s$ es:
\begin{equation}
  \delta_k(\Part_k, s) = \EMD_{\mathrm{efecto}}\!\Bigl(
    \underbrace{p\bigl(\V_{t+1} \mid \V_t = s\bigr)}_{u},\;
    \underbrace{\bigotimes_{i=1}^{k} p\bigl(S_{i,t+1} \mid S_{i,t} = s[M_i]\bigr)}_{v}
  \Bigr)
  \label{eq:delta_k}
\end{equation}

donde $\tensor$ denota el producto tensorial y
$\EMD_{\mathrm{efecto}}(u, v) = \sum_j |u[j] - v[j]|$ es la
Earth Mover's Distance analítica para variables independientes.

\subsubsection*{Significado de cada término}

\begin{orangeBox}
\textbf{En palabras simples:}

\begin{itemize}
  \item \textbf{$s$ (estado actual):} vector binario de longitud $n$ que indica el estado
    presente de cada nodo del sistema; por ejemplo $s = (1, 0, 0, \ldots)$ significa que
    solo el primer nodo está activo.

  \item \textbf{$u = p(\V_{t+1} \mid \V_t = s)$:} distribución de probabilidad del
    sistema \emph{completo} en el instante siguiente, condicionada al estado $s$.
    Captura todas las interacciones causales internas: cada nodo puede influir sobre
    cualquier otro.

  \item \textbf{$v = \bigotimes_{i=1}^{k} p(S_{i,t+1} \mid S_{i,t} = s[M_i])$:}
    distribución del sistema \emph{después de particionarlo} en $k$ partes independientes.
    Cada término $p(S_{i,t+1} \mid S_{i,t} = s[M_i])$ es la distribución de la parte
    $i$ evaluada solo sobre sus propias variables presentes.
    \begin{itemize}
      \item $s[M_i]$ denota el sub-vector de $s$ restringido a los índices del mecanismo
        $M_i$ de la parte $i$; es decir, la parte $i$ ``solo mira'' los nodos que le
        corresponden en el estado actual, ignorando el resto del sistema.
      \item El operador $\otimes$ (producto tensorial) combina las $k$ distribuciones
        marginales como si fueran completamente independientes entre sí.
    \end{itemize}

  \item \textbf{$\EMD(u,v) = \sum_j |u[j] - v[j]|$:} suma de diferencias absolutas
    entre ambas distribuciones, evaluada posición a posición sobre los $2^n$ estados
    posibles del sistema.
\end{itemize}

\medskip
\textbf{La intuición:}
\begin{itemize}
  \item Si la partición es \textbf{buena} (el sistema realmente opera como $k$ partes
    independientes), entonces $u \approx v$ y $\delta_k \approx 0$.
  \item Si la partición es \textbf{mala} (el sistema está integrado causalmente),
    las distribuciones $u$ y $v$ difieren mucho y $\delta_k$ es alta.
\end{itemize}

$\delta_k$ mide entonces \textbf{cuánta información causal se destruye} al asumir que el
sistema opera como $k$ partes independientes en lugar de como un todo integrado.
El k-MIP es la partición que destruye la menor cantidad posible.
\end{orangeBox}

\begin{remark}
Para $k = 2$, la Ecuación~\eqref{eq:delta_k} se reduce exactamente a la función de
pérdida original de GeoMIP para bi-particiones.
\end{remark}

\subsection{Números de Stirling del segundo tipo}

El número de k-particiones de $n$ elementos está dado por el
\textbf{número de Stirling del segundo tipo} $\Stir{n}{k}$:
\begin{equation}
  \Stir{n}{k} = k \cdot \Stir{n-1}{k} + \Stir{n-1}{k-1}, \qquad
  \Stir{n}{1} = \Stir{n}{n} = 1
  \label{eq:stirling}
\end{equation}

\begin{table}[ht]
\centering
\caption{Valores de $\Stir{n}{k}$ y escalabilidad del espacio de búsqueda}
\label{tab:stirling}
\begin{tabular}{@{}r r r r r@{}}
\toprule
$n$ & $k=2$ & $k=3$ & $k=4$ & $k=5$ \\
\midrule
3  & 3      & 1       & ---       & ---       \\
4  & 7      & 6       & 1         & ---       \\
5  & 15     & 25      & 10        & 1         \\
8  & 127    & 966     & 1\,701    & 1\,050    \\
10 & 511    & 9\,330  & 34\,105   & 42\,525   \\
15 & 16\,383 & 2\,375\,101 & \emph{enorme} & \emph{enorme} \\
\bottomrule
\end{tabular}
\end{table}

La cardinalidad total del espacio de búsqueda del subsistema es
$\Stir{|A|}{k} \times \Stir{|M|}{k}$. Para $n=10$, $k=3$:
$9\,330^2 \approx 87$ millones $\Rightarrow$ búsqueda heurística imprescindible.

\subsection{Interpretación geométrica de k-particiones}

\begin{proposition}
En el espacio de estados representado como hipercubo $n$-dimensional $\mathcal{H}^n$,
una k-partición corresponde a la división del hipercubo mediante $k-1$ hiperplanos,
creando $k$ regiones disjuntas.
\end{proposition}

\subsubsection*{El hipercubo de estados}

Cada estado del sistema es un vértice del hipercubo $n$-dimensional
$\mathcal{H}^n = \{0,1\}^n$. Para $n = 3$ nodos (A, B, C), el espacio de estados
tiene $2^3 = 8$ vértices: $(0,0,0), (1,0,0), (0,1,0), \ldots, (1,1,1)$.
La TPM asigna a cada vértice una distribución de probabilidad sobre los $2^n$ vértices
del instante siguiente.

\subsubsection*{Cómo una k-partición divide el hipercubo}

Una k-partición $\Part_k = \{(A_1,M_1), \ldots, (A_k,M_k)\}$ divide el hipercubo en
dos sentidos simultáneos:

\begin{enumerate}
  \item \textbf{Dimensiones futuras (ejes de alcance):}
    El conjunto $A_i$ determina qué \emph{ejes del hipercubo futuro} pertenecen a la
    parte $i$. En términos geométricos, la parte $i$ opera sobre una \emph{cara} del
    hipercubo de dimensión $|A_i|$.

  \item \textbf{Dimensiones presentes (ejes de mecanismo):}
    El conjunto $M_i$ determina qué \emph{ejes del hipercubo presente} puede observar
    la parte $i$. Las dimensiones $\V \setminus M_i$ son ``invisibles'' para esa parte:
    al marginalizar sobre ellas, el hipercubo de la parte $i$ se proyecta sobre un
    subespacio de dimensión $|M_i|$.
\end{enumerate}

\subsubsection*{El corte geométrico como marginalización}

La operación $c.\mathrm{marginalizar}(\mathrm{dims} \setminus M_i)$ colapsa el
hipercubo $n$-dimensional de cada n-cubo $c$ a un subhipercubo de dimensión $|M_i|$,
promediando las probabilidades sobre las dimensiones eliminadas.
Geométricamente: si el hipercubo original tiene 8 vértices y se marginalizan 2 dimensiones,
el resultado es un subhipercubo de $2^1 = 2$ vértices --- una arista del cubo original.

\begin{example}
Sistema de 3 nodos, k-partición en dos partes:
\begin{itemize}
  \item Parte 1: $A_1 = \{A, B\}$, $M_1 = \{A, B\}$ --- subhipercubo $2$-dimensional (cuadrado)
  \item Parte 2: $A_2 = \{C\}$, $M_2 = \{C\}$ --- subhipercubo $1$-dimensional (arista)
\end{itemize}
La distribución conjunta del sistema k-partido vive en el producto cartesiano de estos
dos subespacios: $\{0,1\}^2 \times \{0,1\}^1$, que tiene $4 \times 2 = 8$ combinaciones ---
el mismo número que el espacio original, pero organizado como producto de partes independientes.
\end{example}

\subsubsection*{Por qué GeoMIP usa geometría del hipercubo}

El algoritmo GeoMIP (y su extensión K-GeoMIP) aprovecha la estructura geométrica del
hipercubo para calcular eficientemente el vector de costos $\mathbf{c}$:

\begin{itemize}
  \item Cada vértice del hipercubo de $n$ dimensiones corresponde a un estado del sistema.
  \item Un recorrido BFS desde el estado inicial $s_0$ hasta el estado final $s_f$ sobre el
    hipercubo identifica qué \emph{caras} (subconjuntos de dimensiones) contribuyen más
    a la transición.
  \item El costo $c_i$ de la variable $i$ cuantifica cuánto ``peso'' tiene esa dimensión
    en el recorrido BFS: variables con $c_i$ alto son candidatas a separarse en partes
    distintas porque concentran la mayor parte de la transición causal.
\end{itemize}

Esta representación geométrica permite reducir el problema combinatorio de evaluar
$\Stir{n}{k}^2$ particiones a un recorrido sobre $n \cdot 2^n$ vértices del hipercubo,
que es la fuente de la complejidad $\BigTheta{n \cdot 2^n}$ del algoritmo.

% ════════════════════════════════════════════════════════════
%  SECCIÓN 3 — ARQUITECTURA DEL SOFTWARE
% ════════════════════════════════════════════════════════════
\section{Arquitectura del Software}
\label{sec:arquitectura}

\subsection{Diagrama de arquitectura en capas}

\begin{figure}[H]
\centering
\begin{tikzpicture}[
  layer/.style={
    draw, rounded corners=4pt, align=center,
    text width=13.5cm, minimum height=1.8cm, font=\normalsize
  },
  arr/.style={-{Stealth}, thick}
]
\node[layer, fill=KGeoBlue!20, draw=KGeoBlue] (C1) at (0,0)
  {\textbf{Capa 1 -- Entrada y Configuración} \\
   \texttt{exec.py / main.py} \quad \texttt{Manager, Application}};
\node[layer, fill=KGeoBlue!10, draw=KGeoBlue!70, below=0.45cm of C1] (C2)
  {\textbf{Capa 2 -- Contrato Base / SIA} \\
   \texttt{sia.py} \quad \texttt{sia\_preparar\_subsistema(), chequear\_parametros()}};
\node[layer, fill=KGeoGreen!15, draw=KGeoGreen!70, below=0.45cm of C2] (C3)
  {\textbf{Capa 3 -- Núcleo Computacional IIT} \\
   \texttt{system.py, ncube.py} \quad
   \texttt{System.k\_partir()} \textcolor{KGeoOrange}{\textbf{[NUEVO]}}};
\node[layer, fill=KGeoOrange!15, draw=KGeoOrange!70, below=0.45cm of C3] (C4)
  {\textbf{Capa 4 -- Estrategias Computacionales} \\
   \texttt{controllers/strategies/} \quad
   \texttt{KGeometricSIA} \textcolor{KGeoOrange}{\textbf{[NUEVO]}}};
\node[layer, fill=violet!10, draw=violet!60, below=0.45cm of C4] (C5)
  {\textbf{Capa 5 -- Generadores y Evaluación} \\
   \texttt{funcs/system.py, funcs/base.py} \quad
   \texttt{stirling(), k\_particiones()} \textcolor{KGeoOrange}{\textbf{[NUEVO]}}};
\node[layer, fill=gray!10, draw=gray!60, below=0.45cm of C5] (C6)
  {\textbf{Capa 6 -- Salida y Resultados} \\
   \texttt{solution.py, format.py} \quad
   \texttt{Solution, fmt\_k\_particion()} \textcolor{KGeoOrange}{\textbf{[NUEVO]}}};
\draw[arr, KGeoBlue]    (C1.south) -- (C2.north);
\draw[arr, KGeoBlue!70] (C2.south) -- (C3.north);
\draw[arr, KGeoGreen!70](C3.south) -- (C4.north);
\draw[arr, KGeoOrange!70](C4.south) -- (C5.north);
\draw[arr, violet!60]   (C5.south) -- (C6.north);
\end{tikzpicture}
\caption{Arquitectura en seis capas del framework K-GeoMIP.}
\label{fig:capas}
\end{figure}

\subsection{Diagrama de clases UML}

La Figura~\ref{fig:clases} muestra la jerarquía de clases principal. Cada clase se
representa con un cuadro de cabecera coloreada y cuerpo con sus atributos y métodos.

\begin{figure}[H]
\centering
\begin{tikzpicture}[
  hdr/.style={draw, fill=#1, minimum width=5.4cm, minimum height=0.62cm,
              font=\small\bfseries, text centered, rounded corners=2pt},
  bdy/.style={draw, fill=white, minimum width=5.4cm, align=left,
              font=\ttfamily\scriptsize, inner sep=4pt, text width=5.1cm},
  arr/.style={-{Stealth}, thick, dashed, KGeoGray},
  inh/.style={-{Triangle[open]}, thick},
]
% === SIA (centro-arriba) ===
\node[hdr=KGeoBlue!35] (SIA_h) at (0,6.8) {SIA \small(ABC)};
\node[bdy, below=-\pgflinewidth of SIA_h] (SIA_b) {
  +sia\_preparar\_subsistema()\\
  +chequear\_parametros()\\
  +sia\_cargar\_tpm()\\
  \#sia\_subsistema: System\\
  \#sia\_dists\_marginales
};

% === GeometricSIA (izquierda) ===
\node[hdr=KGeoGreen!30] (GEO_h) at (-3.2,3.5) {GeometricSIA};
\node[bdy, below=-\pgflinewidth of GEO_h] (GEO_b) {
  +aplicar\_estrategia()\\
  +find\_mip()\\
  +calcular\_costos\_nivel()\\
  +calcular\_costo()\\
  +identificar\_particiones()
};

% === KGeometricSIA (derecha) ===
\node[hdr=KGeoOrange!40] (KGEO_h) at (3.2,3.5) {KGeometricSIA \textcolor{KGeoOrange}{*}};
\node[bdy, below=-\pgflinewidth of KGEO_h] (KGEO_b) {
  +aplicar\_estrategia(k)\\
  -\_generar\_candidatos\_k()\\
  -\_greedy\_multiway()\\
  -\_grupos\_a\_particion()\\
  -\_particion\_valida()\\
  -\_cache\_subsistema: dict
};

% === System (abajo-izq) ===
\node[hdr=KGeoBlue!20] (SYS_h) at (-3.2,0.2) {System};
\node[bdy, below=-\pgflinewidth of SYS_h] (SYS_b) {
  +condicionar()\\
  +substraer()\\
  +bipartir()\\
  +k\_partir() \textcolor{KGeoOrange}{*}\\
  +distribucion\_marginal()
};

% === NCube (abajo-der) ===
\node[hdr=KGeoBlue!10] (NC_h) at (3.2,0.2) {NCube};
\node[bdy, below=-\pgflinewidth of NC_h] (NC_b) {
  +condicionar()\\
  +marginalizar()\\
  +indice: np.int8\\
  +dims: NDArray\\
  +data: NDArray
};

% --- Herencia (triángulo hueco apuntando al padre) ---
\draw[inh] (GEO_h.north) -- ++(0,0.4) -| (SIA_b.south west);
\draw[inh] (KGEO_h.north) -- ++(0,0.4) -| (SIA_b.south east);
% Ruta por el hueco entre las dos clases (x=0): evita atravesar cualquier bloque
\draw[inh] (KGEO_h.west) -- node[midway,above,font=\tiny]{hereda} (GEO_h.east);

% --- Dependencias (flecha punteada) ---
\draw[arr] (GEO_b.south) -- (SYS_h.north);
\draw[arr] (SYS_b.east)  -- (NC_h.west);
\end{tikzpicture}
\caption{Diagrama de clases UML de K-GeoMIP.
  El símbolo \textcolor{KGeoOrange}{*} indica componentes nuevos o extendidos.}
\label{fig:clases}
\end{figure}

\subsection{Estructura de directorios del repositorio}

\begin{lstlisting}[style=pseudocode,
  caption={Estructura del repositorio K-GeoMIP (Método 2)},
  label=lst:estructura]
K-GeoMIP/
  src/Method2_Dynamic_Programming_Reformulation/
    exec.py                        <- Punto de entrada
    src/
      main.py                      <- iniciar(), run_prueba()
      constants/
        models.py                  <- K_MIN=2, K_MAX=5, KGEOMETRIC_*
      controllers/
        manager.py                 <- Manager.generar_red() (chunks)
        strategies/
          geometric.py             <- GeometricSIA (k=2)
          k_geometric.py           <- KGeometricSIA [NUEVO]
      models/
        base/
          sia.py                   <- Contrato SIA (ABC)
        core/
          system.py                <- System.k_partir() [NUEVO]
          ncube.py                 <- NCube
          solution.py              <- Solution
      funcs/
        system.py                  <- stirling, k_particiones [NUEVO]
        format.py                  <- fmt_k_particion [NUEVO]
        base.py                    <- emd_efecto, ABECEDARY
    tests/
      conftest.py                  <- Fixtures, fix colorama
      test_kgeomip.py              <- 13 tests (grupos A, B, C)
    experiments/
      benchmark_paso6.py           <- Benchmarks experimentales
    data/samples/
      N3A.csv ... N20A.csv         <- TPMs de prueba
\end{lstlisting}

\subsection{Diagrama de secuencia principal}

\begin{figure}[ht]
\centering
\begin{tikzpicture}[
  obj/.style={draw, fill=KGeoBlue!15, minimum width=2.2cm,
              minimum height=0.55cm, font=\small, text centered},
  life/.style={draw, dashed, KGeoGray},
  fwd/.style={-{Stealth}, thick},
  ret/.style={-{Stealth}, dashed, KGeoGray},
  lbl/.style={font=\scriptsize, above},
]
% Objetos
\node[obj] (caller) at (0,0)    {Llamador};
\node[obj] (kgeo)   at (3.5,0)  {KGeoSIA};
\node[obj] (sys)    at (7.0,0)  {System};
\node[obj] (emd)    at (10.5,0) {emd\_efecto};
% Líneas de vida
\draw[life] (caller.south) -- ++(0,-10);
\draw[life] (kgeo.south)   -- ++(0,-10);
\draw[life] (sys.south)    -- ++(0,-10);
\draw[life] (emd.south)    -- ++(0,-10);
% --- mensajes ---
\draw[fwd]($(caller.south)+(0,-0.5)$)--
  node[lbl]{aplicar\_estrategia(k>2)}($(kgeo.south)+(0,-0.5)$);
\draw[fwd]($(kgeo.south)+(0,-1.2)$)--
  node[lbl]{preparar\_subsistema()}($(sys.south)+(0,-1.2)$);
\draw[ret]($(sys.south)+(0,-1.8)$)--
  node[lbl]{System(sub)}($(kgeo.south)+(0,-1.8)$);

\node[font=\scriptsize\itshape,KGeoGray] at (-0.8,-2.4) {[Fase 3]};
\draw[fwd]($(kgeo.south)+(0,-2.6)$)--
  node[lbl]{calcular\_costos\_nivel()}($(sys.south)+(0,-2.6)$);
\draw[ret]($(sys.south)+(0,-3.2)$)--
  node[lbl]{\_trans\_matrix}($(kgeo.south)+(0,-3.2)$);

\node[font=\scriptsize\itshape,KGeoGray] at (-0.8,-3.8) {[Fase 4]};
\draw[fwd]($(kgeo.south)+(0,-4.0)$)--
  node[lbl]{\_generar\_candidatos\_k(k)}($(sys.south)+(0,-4.0)$);
\draw[ret]($(sys.south)+(0,-4.6)$)--
  node[lbl]{[p1,...,p25]}($(kgeo.south)+(0,-4.6)$);

\node[font=\scriptsize\itshape,KGeoGray] at (-0.8,-5.2) {[Fase 5]};
\draw[fwd]($(kgeo.south)+(0,-5.4)$)--
  node[lbl]{k\_partir(pi)}($(sys.south)+(0,-5.4)$);
\draw[ret]($(sys.south)+(0,-6.0)$)--
  node[lbl]{System(k-partido)}($(kgeo.south)+(0,-6.0)$);
\draw[fwd]($(kgeo.south)+(0,-6.6)$)--
  node[lbl]{emd\_efecto(dist, orig)}($(emd.south)+(0,-6.6)$);
\draw[ret]($(emd.south)+(0,-7.2)$)--
  node[lbl]{$\delta_k$}($(kgeo.south)+(0,-7.2)$);

\draw[ret]($(kgeo.south)+(0,-8.6)$)--
  node[lbl]{Solution($\delta_{\min}$, k-MIP)}($(caller.south)+(0,-8.6)$);
\end{tikzpicture}
\caption{Diagrama de secuencia de \texttt{KGeometricSIA.aplicar\_estrategia()} para $k > 2$.}
\label{fig:secuencia}
\end{figure}

\subsection{Patrones de Diseño Aplicados}

\begin{table}[H]
\centering
\caption{Patrones de diseño utilizados en K-GeoMIP}
\label{tab:patrones}
\begin{tabular}{@{}l p{4.2cm} p{5cm}@{}}
\toprule
\textbf{Patrón} & \textbf{Dónde se aplica} & \textbf{Beneficio} \\
\midrule
\textbf{Strategy} &
  \texttt{GeometricSIA}, \texttt{KGeometricSIA} heredan de \texttt{SIA} y
  proveen su propia \texttt{aplicar\_estrategia()} &
  Intercambiabilidad sin modificar \texttt{Manager} \\[4pt]
\textbf{Template Method} &
  \texttt{SIA} define el esqueleto; subclases rellenan \texttt{aplicar\_estrategia()} &
  Reutilización del flujo de preparación del subsistema \\[4pt]
\textbf{Factory Method} &
  \texttt{Manager} instancia la estrategia según \texttt{KGEOMIP\_K} &
  Desacoplamiento entre config.\ y clase concreta \\[4pt]
\textbf{Iterator/Generator} &
  \texttt{k\_particiones()} usa \texttt{yield} &
  Carga perezosa: $O(n)$ RAM independiente de $\Stir{n}{k}$ \\[4pt]
\textbf{Flyweight/Cache} &
  \texttt{\_cache\_subsistema} en \texttt{KGeometricSIA} &
  Elimina recómputo del BFS entre distintos $k$ (ahorro 75\%) \\
\bottomrule
\end{tabular}
\end{table}

\subsection{Decisiones Arquitectónicas Clave}

\subsubsection{Reutilización de la tabla BFS geométrica}
\textbf{Decisión:} \texttt{KGeometricSIA} llama a \texttt{calcular\_costos\_nivel()} de
\texttt{GeometricSIA} sin reimplementarlo.

\textbf{Motivación:} La tabla de costos $\mathbf{C}[s_0, s_f]$ es independiente de $k$;
captura el peso de cada variable respecto al estado inicial, que es el mismo insumo para
la heurística LPT sin importar cuántas partes se definan.

\textbf{Trade-off:} KGeometricSIA queda acoplado a la implementación interna de
GeometricSIA. Se aceptó dado que GeometricSIA es la implementación de referencia estable.

\subsubsection{Generadores perezosos vs.\ listas materializadas}
\textbf{Decisión:} \texttt{k\_particiones()} es un generador Python.

\textbf{Motivación:} Para $n=10$, $k=3$, materializar todas las particiones requeriría
$9330^2 \times 10 \times 4\,\text{bytes} \approx 3{,}5\,\text{GB}$ de RAM.

\subsubsection{$k=2$ delegado a GeometricSIA}
\textbf{Decisión:} Si el llamador pasa $k=2$, \texttt{KGeometricSIA} delega
inmediatamente a \texttt{GeometricSIA.\allowbreak{}aplicar\_estrategia()},
que hace búsqueda \textbf{exacta} (no heurística) sobre bi-particiones.

\subsubsection{Constante $C = 25$ candidatos}
\textbf{Decisión:} Número total de candidatos fijo en $C = 25$
(3 deterministas + hasta 22 perturbaciones aleatorias con semilla reproducible).

\textbf{Motivación:} Para los datasets de prueba, 25 evaluaciones de EMD toman
$< 2\,\text{ms}$, mientras que el BFS toma $\sim 100\,\text{ms}$ ($n=10$). Aumentar
$C$ a 100 solo agregaría $\sim 7\,\text{ms}$ de impacto mínimo.

% ════════════════════════════════════════════════════════════
%  SECCIÓN 4 — DISEÑO ALGORÍTMICO
% ════════════════════════════════════════════════════════════
\section{Diseño Algorítmico}
\label{sec:diseno}

\subsection{\texorpdfstring{\texttt{System.k\_partir()}}{System.k\_partir()} --- Núcleo}

Sea un subsistema con n-cubos $\mathbf{c} = (c_0, \ldots, c_{n-1})$. Dado
$\mathrm{particion} = [(A_1, M_1), \ldots, (A_k, M_k)]$, la transformación es:

\begin{equation}
  \forall\, c_j: \quad c_j' \leftarrow c_j.\mathrm{marginalizar}
  \bigl(c_j.\mathrm{dims} \setminus M_{i^*}\bigr), \quad
  i^* = \text{único } i \text{ tal que } c_j.\mathrm{índice} \in A_i
  \label{eq:k_partir}
\end{equation}

\begin{proposition}
\texttt{k\_partir}$\bigl([(A_1, M_1), (\V \setminus A_1, \V \setminus M_1)]\bigr)
\equiv \texttt{bipartir}(A_1, M_1)$.
\end{proposition}

\begin{algorithm}[H]
\SetAlgoLined
\KwIn{$\mathrm{particion} = [(A_1, M_1), \ldots, (A_k, M_k)]$}
\KwOut{Nuevo \texttt{System} con n-cubos marginalizados}
$\mathrm{new\_sys} \leftarrow \texttt{System.\_\_new\_\_}(\texttt{System})$\;
$\mathrm{new\_sys.estado\_inicial} \leftarrow \mathrm{self.estado\_inicial}$\;
$\mathrm{indice\_a\_mec} \leftarrow \{\}$\;
\ForEach{$(A_i, M_i) \in \mathrm{particion}$}{
  \ForEach{$\mathrm{idx} \in A_i$}{
    $\mathrm{indice\_a\_mec}[\mathrm{int(idx)}] \leftarrow M_i$\;
  }
}
$\mathrm{new\_sys.ncubos} \leftarrow \texttt{tuple}($
  $\mathrm{cube.marginalizar}(\mathrm{cube.dims} \setminus
    \mathrm{indice\_a\_mec}[\mathrm{cube.indice}])$
  $\text{ for cube in self.ncubos})$\;
\KwRet{$\mathrm{new\_sys}$}\;
\caption{\texttt{System.k\_partir(particion)}}
\label{alg:k_partir}
\end{algorithm}

\subsection{Generadores de k-particiones}

\subsubsection{\texttt{stirling(n, k)} --- Programación dinámica}

\begin{algorithm}[H]
\SetAlgoLined
\KwIn{$n$ (tamaño), $k$ (número de partes)}
\KwOut{$\Stir{n}{k} \in \N_0$}
\If{$k = 0$ \textbf{or} $k > n$}{\KwRet{$0$}}
\If{$k = 1$ \textbf{or} $k = n$}{\KwRet{$1$}}
$\mathrm{prev}[0..k] \leftarrow [0, 1, 0, \ldots, 0]$\;
\For{$i \leftarrow 2$ \KwTo $n$}{
  $\mathrm{curr}[j] \leftarrow j \cdot \mathrm{prev}[j] + \mathrm{prev}[j-1]
    \quad \forall\, j \in [1, \min(i,k)]$\;
  $\mathrm{prev} \leftarrow \mathrm{curr}$\;
}
\KwRet{$\mathrm{prev}[k]$}\;
\caption{\texttt{stirling(n, k)} --- O(n·k) tiempo, O(k) espacio}
\label{alg:stirling}
\end{algorithm}

\subsubsection{\texttt{particionar\_conjunto(elementos, k)} --- Generador RGS}

\paragraph{Definición formal.}
Una \textbf{Cadena de Crecimiento Restringido} (\emph{Restricted Growth String}, RGS)
es un arreglo $\mathbf{a} = [a_0, a_1, \ldots, a_{n-1}]$ donde $a_i \in \{0,\ldots,k-1\}$
indica a qué grupo pertenece el $i$-ésimo elemento, sujeto a dos invariantes:

\begin{center}
\begin{tabular}{@{}l p{9cm}@{}}
\toprule
\textbf{Invariante} & \textbf{Por qué} \\
\midrule
$a_0 = 0$ siempre &
  El primer elemento siempre se asigna al grupo~0.
  Esto \emph{rompe la simetría de los grupos}: sin esta regla,
  $\{0{,}1\},\{2\}$ y $\{2\},\{0{,}1\}$ serían dos RGS distintas
  para la misma partición lógica. \\[4pt]
$a_i \leq \max(a_0,\ldots,a_{i-1}) + 1$ &
  No se puede ``saltar'' un número de grupo. Si el grupo~1 no ha
  aparecido antes, no se puede asignar el grupo~2: eso dejaría
  el grupo~1 vacío, violando la no-trivialidad
  (condición~\ref{eq:notrivial}). \\
\bottomrule
\end{tabular}
\end{center}

La RGS es \textbf{válida para $k$ partes} cuando $\max(\mathbf{a}) = k-1$
(todos los grupos $0, 1, \ldots, k-1$ han sido usados al menos una vez).

\begin{proposition}[Biyección RGS-Particiones]
Existe una correspondencia biunívoca entre las RGS con $\max(\mathbf{a}) = k-1$
y las k-particiones de $\{0, \ldots, n-1\}$.
El número de RGS válidas es exactamente $\Stir{n}{k}$.
\end{proposition}

\begin{orangeBox}
\textbf{¿Qué significa ``correspondencia biunívoca''?}

Una correspondencia \textbf{biunívoca} (o \emph{biyección}) entre dos conjuntos $X$ e $Y$
es una función $f: X \to Y$ que satisface simultáneamente dos propiedades:

\begin{itemize}
  \item \textbf{Inyectiva} (\emph{uno a uno}): elementos distintos de $X$ producen
    imágenes distintas en $Y$.
    $x_1 \neq x_2 \;\Rightarrow\; f(x_1) \neq f(x_2)$.
    En nuestro caso: dos RGS distintas $\mathbf{a} \neq \mathbf{b}$ generan
    particiones distintas. \textbf{Sin duplicados.}

  \item \textbf{Sobreyectiva} (\emph{sobre}): todo elemento de $Y$ tiene al menos
    una preimagen en $X$.
    $\forall\, y \in Y,\; \exists\, x \in X : f(x) = y$.
    En nuestro caso: toda k-partición posible de $\{0,\ldots,n-1\}$ tiene exactamente
    una RGS canónica que la representa. \textbf{Sin omisiones.}
\end{itemize}

\medskip
\textbf{Por qué se cumple en RGS:}
\begin{enumerate}
  \item \textbf{Inyectividad}: la regla $a_0 = 0$ fija un representante canónico
    por cada clase de equivalencia de grupos (elimina permutaciones de etiquetas).
    La regla de crecimiento $a_i \leq \max_{j<i}(a_j) + 1$ garantiza que no haya
    dos cadenas distintas que describan la misma asignación de elementos a grupos.

  \item \textbf{Sobreyectividad}: dada cualquier k-partición
    $\{S_0, S_1, \ldots, S_{k-1}\}$, basta ordenar sus partes de modo que
    $\min(S_0) < \min(S_1) < \cdots < \min(S_{k-1})$ (el grupo~$i$ es el que
    contiene el $i$-ésimo mínimo), y asignar $a[j] = i$ si el elemento~$j$
    pertenece a $S_i$. El resultado es siempre una RGS válida con $\max = k-1$.
\end{enumerate}

\medskip
\textbf{Consecuencia práctica:} el DFS sobre RGS genera \emph{exactamente}
$\Stir{n}{k}$ cadenas válidas --- el mismo número que k-particiones existen ---
sin repetir ninguna ni omitir ninguna. Esto hace que el generador sea correcto
por construcción, sin necesidad de desduplicar a posteriori.
\end{orangeBox}

\paragraph{Ejemplo trazado: elementos $= \{0,1,2\}$, $k = 2$.}

\begin{blueBox}
\textbf{Problema:} Generar todas las 2-particiones de $\{0,1,2\}$.
Resultado esperado: $\Stir{3}{2} = 3$ particiones.

\medskip
Las tres particiones y sus RGS canónicas:
\begin{center}
\begin{tabular}{@{}l l@{}}
\toprule
\textbf{Partición} & \textbf{RGS} $\mathbf{a}$ \\
\midrule
$S_1=\{0,1\},\; S_2=\{2\}$ & $[0, 0, 1]$ \\
$S_1=\{0,2\},\; S_2=\{1\}$ & $[0, 1, 0]$ \\
$S_1=\{0\},\;   S_2=\{1,2\}$ & $[0, 1, 1]$ \\
\bottomrule
\end{tabular}
\end{center}
\end{blueBox}

\begin{lstlisting}[style=pseudocode,
  caption={Árbol de recursión de \texttt{\_rgs([0,1,2], k=2)} --- 3 hojas válidas},
  label=lst:rgs-trace]
a = [0, ?, ?]         <- a[0]=0 fijo; llamada inicial: _rgs(pos=1, max=0)

_rgs(pos=1, max=0)    v puede ser 0..min(max+1,k-1) = 0..1
  |
  +-- v=0: a=[0,0,?]  _rgs(pos=2, max=0)
  |     |
  |     +-- v=0: a=[0,0,0]  max=0 != k-1=1  INVALIDO (falta grupo 1)
  |     |
  |     +-- v=1: a=[0,0,1]  max=1 == k-1=1  ** VALIDO **
  |               Parte_0={0,1}  Parte_1={2}
  |
  +-- v=1: a=[0,1,?]  _rgs(pos=2, max=1)
        |
        +-- v=0: a=[0,1,0]  max=1 == k-1=1  ** VALIDO **
        |         Parte_0={0,2}  Parte_1={1}
        |
        +-- v=1: a=[0,1,1]  max=1 == k-1=1  ** VALIDO **
                  Parte_0={0}    Parte_1={1,2}

Total: 3 particiones = S(3,2)  (sin repeticiones, sin omisiones)
\end{lstlisting}

\paragraph{Cómo se traduce la RGS a una partición.}
Cada posición $i$ del arreglo $\mathbf{a}$ indica directamente a qué parte va
el elemento $i$. Por ejemplo, $\mathbf{a} = [0, 0, 1]$:
\begin{itemize}
  \item $a[0] = 0 \;\Rightarrow$ elemento~$0 \to$ Parte$_0$
  \item $a[1] = 0 \;\Rightarrow$ elemento~$1 \to$ Parte$_0$
  \item $a[2] = 1 \;\Rightarrow$ elemento~$2 \to$ Parte$_1$
\end{itemize}
Resultado: Parte$_0 = \{0,1\}$, Parte$_1 = \{2\}$
$\;\longrightarrow\;$ partición $\bigl(\{0,1\},\{2\}\bigr)$.

Gracias a que $a_0 = 0$ es fijo y la regla de crecimiento impone un orden canónico
sobre los grupos, \emph{cada partición tiene exactamente una RGS} (sin duplicados) y
\emph{cada RGS produce exactamente una partición} (sin omisiones). El DFS enumera todas
las RGS válidas con $\max = k-1$ de forma perezosa vía \texttt{yield}, consumiendo
solo $\BigO{n}$ de memoria en la pila independientemente de $\Stir{n}{k}$.

\begin{algorithm}[H]
\SetAlgoLined
\KwIn{$\mathrm{elementos}$, $k$}
\KwOut{Generador de k-particiones (yield)}
$\mathbf{a}[0..n-1] \leftarrow \mathbf{0}$\;
\SetKwFunction{FRGS}{rgs}
\SetKwProg{Fn}{Función}{:}{}
\Fn{\FRGS{pos, max\_usado}}{
  \If{pos $=$ n}{
    \If{max\_usado $=$ k-1}{
      $\mathrm{partes} \leftarrow [[\,],\ldots,[\,]]$\;
      yield $\texttt{tuple}(\mathrm{partes})$\;
    }
    \KwRet{}
  }
  \For{$v \leftarrow 0$ \KwTo $\min(\mathrm{max\_usado}+1,\; k-1)$}{
    $a[\mathrm{pos}] \leftarrow v$\;
    \FRGS{pos+1, $\max(\mathrm{max\_usado}, v)$}\;
  }
}
\FRGS{1, 0}\;
\caption{\texttt{particionar\_conjunto(elementos, k)} --- DFS sobre RGS, $\BigO{n}$ espacio}
\label{alg:rgs}
\end{algorithm}

\subsection{Estrategia heurística: \texttt{KGeometricSIA}}
\label{sec:kgeomip}

\subsubsection{Principio del algoritmo LPT}

\paragraph{Definición.}
\textbf{LPT} (\emph{Longest Processing Time first}) es una heurística greedy para
repartir $n$ variables en $k$ grupos minimizando el \textbf{makespan}
(peso del grupo más cargado):
\begin{equation}
  \min\; \max_{j=0}^{k-1} \sum_{i \in G_j} c_i
  \label{eq:lpt-objetivo}
\end{equation}

\begin{blueBox}
\textbf{Analogía:} Imagina $k$ mochilas vacías y $n$ objetos con pesos distintos.
LPT toma primero el objeto \emph{más pesado} y lo coloca en la mochila con menor
peso acumulado. Luego el segundo más pesado a la más vacía, y así sucesivamente.
Al terminar, las mochilas tienden a tener pesos similares --- el desbalance es mínimo.
\end{blueBox}

\paragraph{Algoritmo paso a paso.}
Dado el vector de costos $\mathbf{c}$ y el número de grupos $k$:
\begin{enumerate}
  \item \textbf{Ordenar} las variables de \textbf{mayor a menor} costo:
    $c_{\sigma(0)} \geq c_{\sigma(1)} \geq \cdots \geq c_{\sigma(n-1)}$.
  \item \textbf{Mantener} un \emph{heap mínimo} sobre el peso acumulado de los $k$ grupos
    (inicialmente todos en cero).
  \item \textbf{Para cada variable} en orden descendente:
    extraer el grupo de menor peso del heap, asignarle la variable,
    actualizar su peso y reinsertar en el heap.
\end{enumerate}
Complejidad total: $\BigO{n \log n}$ (dominada por el ordenamiento inicial).

\paragraph{Ejemplo concreto --- variables $[a,b,c,d,e]$, costos $[5,3,4,2,1]$, $k=3$.}

\textbf{Paso 1 --- Ordenar desc:} $a(5),\; c(4),\; b(3),\; d(2),\; e(1)$.

\textbf{Paso 2 --- Asignación greedy} con heap mínimo sobre sumas de grupos:

\begin{center}
\begin{tabular}{@{}c l r r r r p{4.2cm}@{}}
\toprule
\textbf{It.} & \textbf{Var.} & \textbf{Costo} &
$\Sigma G_0$ & $\Sigma G_1$ & $\Sigma G_2$ & \textbf{Acción} \\
\midrule
1 & $a$ & 5 & \textbf{5} & 0 & 0 &
  heap min $= 0$ ($G_0$) $\to$ asignar a $G_0$ \\
2 & $c$ & 4 & 5 & \textbf{4} & 0 &
  heap min $= 0$ ($G_1$) $\to$ asignar a $G_1$ \\
3 & $b$ & 3 & 5 & 4 & \textbf{3} &
  heap min $= 0$ ($G_2$) $\to$ asignar a $G_2$ \\
4 & $d$ & 2 & 5 & 4 & \textbf{5} &
  heap min $= 3$ ($G_2$) $\to$ asignar a $G_2$ \\
5 & $e$ & 1 & 5 & \textbf{5} & 5 &
  heap min $= 4$ ($G_1$) $\to$ asignar a $G_1$ \\
\bottomrule
\end{tabular}
\end{center}

\textbf{Resultado:}
$G_0 = \{a\} = 5$,\quad
$G_1 = \{c,\, e\} = 5$,\quad
$G_2 = \{b,\, d\} = 5$
$\;\Rightarrow\;$ makespan $= \max(5,5,5) = 5$ \quad\emph{(balance perfecto en este caso)}.

\paragraph{Garantía formal de Graham (1969).}

\begin{blueBox}
\begin{equation}
  \mathrm{makespan}(\mathrm{LPT}) \leq
  \left(\frac{4}{3} - \frac{1}{3k}\right) \cdot \mathrm{OPT}
  \label{eq:lpt-garantia}
\end{equation}
\begin{itemize}
  \item $k=3$: makespan $\leq 1{,}22 \cdot \mathrm{OPT}$
    (a lo sumo 22\% peor que el óptimo).
  \item $k=5$: makespan $\leq 1{,}27 \cdot \mathrm{OPT}$
    (a lo sumo 27\% peor).
\end{itemize}
Esta es la cota más ajustada conocida para LPT: demostrable que no existe
garantía mejor en el caso general de $k$ máquinas idénticas.
\end{blueBox}

\paragraph{Advertencia: la garantía LPT es sobre makespan, no sobre EMD.}

\begin{orangeBox}
La cota de Graham aplica sobre el \textbf{balance de costos} (makespan),
\emph{no} sobre la pérdida de información causal $\delta_k$~(EMD):

\begin{itemize}
  \item LPT asegura que los grupos tienen pesos $c_i$ \emph{similares entre sí}.
  \item LPT \textbf{no} garantiza que esa partición sea la de menor $\delta_k$,
    porque el EMD captura \emph{correlaciones conjuntas} entre variables,
    mientras que $c_i$ solo refleja el peso marginal de cada variable de forma
    independiente.
\end{itemize}

\medskip
Por eso K-GeoMIP genera \textbf{tres candidatos complementarios}
(LPT, round-robin, bloques) más hasta 22 perturbaciones aleatorias,
y elige entre todos el que minimiza el EMD real evaluando $\delta_k$
directamente sobre cada candidato.
\end{orangeBox}

\subsubsection{Las tres estrategias de generación de candidatos}

\begin{table}[H]
\centering
\caption{Las tres estrategias deterministas de \texttt{\_generar\_candidatos\_k}}
\label{tab:estrategias}
\begin{tabular}{@{}l l l l@{}}
\toprule
\textbf{Estrategia} & \textbf{Descripción} & \textbf{Costo} & \textbf{Garantía} \\
\midrule
$E_1$ --- LPT & Sort $\downarrow$, heap mínimo & $\BigO{n \log n}$ &
  makespan $\leq (4/3 - 1/(3k))\cdot\mathrm{OPT}$ \\
$E_2$ --- Round-robin & Sort $\uparrow$, asignación cíclica & $\BigO{n \log n}$ &
  Distribución uniforme \\
$E_3$ --- Bloques & Bloques de $\lfloor n/k \rfloor$ por índice & $\BigO{n}$ &
  Dependencias contiguas \\
\bottomrule
\end{tabular}
\end{table}

Más hasta $C - 3 = 22$ \textbf{perturbaciones aleatorias} con semilla reproducible
$c_i' = c_i \cdot \mathcal{U}(0{,}6;\;1{,}4)$, con deduplicación por firma
\texttt{frozenset}.

\paragraph{Mecanismo adaptativo de $n_{\mathrm{obj}}$.}

El límite de candidatos está definido como constante
\texttt{\_C\_CANDIDATOS\_TOTAL = 25} (línea 37 de \texttt{k\_geometric.py}),
pero el número real de candidatos generados es \textbf{adaptativo}:
\begin{equation}
  n_{\mathrm{obj}} = \min\!\bigl(C,\; \Stir{n}{k}\bigr)
  \label{eq:n_obj}
\end{equation}

Esto significa que si el espacio de particiones es más pequeño que $C = 25$
(por ejemplo, $\Stir{3}{2} = 3$), el generador produce solo las 3 particiones
realmente existentes y no desperdicia iteraciones buscando duplicados.
Por el contrario, si $\Stir{n}{k} \gg 25$ (como $\Stir{10}{3} \approx 9\,330$),
el generador produce los 25 candidatos completos.

\begin{center}
\begin{tabular}{@{}r r r r@{}}
\toprule
$n$ & $k$ & $\Stir{n}{k}$ & $n_{\mathrm{obj}}$ efectivo \\
\midrule
3 & 2 & 3       & 3  \quad {\small(exhaustivo)} \\
4 & 3 & 6       & 6  \quad {\small(exhaustivo)} \\
5 & 3 & 25      & 25 \quad {\small(exactamente C)} \\
6 & 3 & 90      & 25 \quad {\small(heurístico)} \\
10 & 3 & 9\,330 & 25 \quad {\small(heurístico)} \\
\bottomrule
\end{tabular}
\end{center}

\paragraph{¿Por qué $C = 25$ y no más?}

\begin{blueBox}
El tiempo dominante del algoritmo \textbf{no} son las perturbaciones
($\BigO{n \log n}$ por candidato), sino la \textbf{evaluación EMD}
de cada candidato ($\BigO{n \cdot 2^n}$ por candidato):

\begin{itemize}
  \item Para $n=10$: evaluar 1 candidato cuesta $\approx 10 \cdot 1024$ operaciones.
    Evaluar 25 cuesta $25 \times$ eso --- aun así, esto es apenas una fracción del
    costo del BFS ($\BigTheta{n \cdot 2^n}$ fijo).
  \item Aumentar $C$ de 25 a 100 multiplicaría el tiempo de evaluación de candidatos
    por $4\times$, pero con beneficio marginal casi nulo en calidad: las perturbaciones
    con $\mathcal{U}(0{,}6;\;1{,}4)$ ya cubren suficiente variedad del espacio de
    búsqueda para $n \leq 12$.
  \item El speedup versus búsqueda exhaustiva se mantiene en
    $\Stir{n}{k}^2 / C \approx 3{,}47 \times 10^6$ para $n=10$, $k=3$
    (ver Ecuación~\ref{eq:speedup}), independientemente de si $C = 25$ o $C = 100$.
\end{itemize}

\medskip
$C = 25$ es el punto de equilibrio \textbf{costo-beneficio}: suficiente variedad
de candidatos para obtener una buena aproximación, sin que el costo de evaluación
EMD supere al del BFS que domina la complejidad total.
\end{blueBox}

\subsubsection{Pseudocódigo completo de \texttt{aplicar\_estrategia}}

\begin{algorithm}[H]
\SetAlgoLined
\KwIn{condicion, alcance, mecanismo, tpm, $k \in \{2,\ldots,5\}$}
\KwOut{\texttt{Solution} con $\delta_k$ mínima}
\If{$k = 2$}{\KwRet{\texttt{GeometricSIA.aplicar\_estrategia}(...)}\;}
\tcc{Fase 1: preparar subsistema}
\texttt{sia\_preparar\_subsistema(condicion, alcance, mecanismo, tpm)}\;
\tcc{Fase 3: tabla de costos BFS --- $\BigTheta{n \cdot 2^n}$}
\For{nivel $\leftarrow 1$ \KwTo $n$}{
  \texttt{calcular\_costos\_nivel}(estado\_final, nivel)\;
}
\tcc{Fase 4: generar $C = 25$ candidatos}
$\mathbf{c} \leftarrow \texttt{\_trans\_matrix}[s_0, s_f]$\;
candidatos $\leftarrow [\texttt{LPT}(\mathbf{c}),\; \texttt{round\_robin}(\mathbf{c}),\;
  \texttt{bloques}(n,k)] + \text{perturbaciones}$\;
\tcc{Fase 5: evaluar EMD para cada candidato}
\ForEach{$P \in \mathrm{candidatos}$}{
  $\delta_k(P) \leftarrow \texttt{emd\_efecto}(
    \texttt{k\_partir}(P).\texttt{dist\_marginal}(),\; \mathrm{dist\_sub})$\;
}
\tcc{Fase 6: seleccionar k-MIP}
$P^* \leftarrow \argmin_{P} \delta_k(P)$\;
\KwRet{\texttt{Solution}(``KGeometric'', $\delta_k(P^*)$,
  \texttt{fmt\_k\_particion}($P^*$), $t$)}\;
\caption{\texttt{KGeometricSIA.aplicar\_estrategia}}
\label{alg:kgeomip}
\end{algorithm}

% ════════════════════════════════════════════════════════════
%  SECCIÓN 5 — ANÁLISIS DE COMPLEJIDAD
% ════════════════════════════════════════════════════════════
\section{Análisis de Complejidad}
\label{sec:complejidad}

\subsection{Complejidad temporal por fase}

\begin{table}[ht]
\centering
\caption{Análisis de complejidad por fase de \texttt{KGeometricSIA.aplicar\_estrategia}}
\label{tab:complejidad-fases}
\begin{tabular}{@{}l l l l@{}}
\toprule
\textbf{Fase} & \textbf{Operación} & \textbf{$T(n)$} & \textbf{Notación} \\
\midrule
1 -- Preparar subsistema & condicionar + substraer & $\BigTheta{n \cdot 2^n}$ & $\Theta$ \\
3 -- Tabla de costos BFS & calcular\_costos\_nivel $\times n$ & $\BigTheta{n \cdot 2^n}$ & $\Theta$ \\
4 -- Generar candidatos  & Sort + heap $\times C$ & $\BigO{C \cdot n \log n}$ & $O$ \\
5 -- Evaluar EMD $\times C$ & k\_partir + emd\_efecto & $\BigO{C \cdot n \cdot 2^n}$ & $O$ \\
6 -- Seleccionar mínimo & min sobre $C$ claves & $\BigO{1}$ & $O$ \\
\midrule
\textbf{Total} & & $\bm{\BigTheta{n \cdot 2^n}}$ & $\bm{\Theta}$ \\
\bottomrule
\end{tabular}
\end{table}

\subsection{Demostración de la cota $\BigTheta{n \cdot 2^n}$}

\begin{align}
  T(n) &= T_3(n) + T_4(n) + T_5(n) \notag \\
       &= \BigTheta{n \cdot 2^n} + \BigO{n \log n} + \BigO{C \cdot n \cdot 2^n}
  \label{eq:complejidad-total}
\end{align}

Como $n \log n = o(n \cdot 2^n)$ y $C = 25$ es constante:
\begin{equation}
  \exists\, c_1, c_2 > 0: \quad
  c_1 \cdot n \cdot 2^n \leq T(n) \leq c_2 \cdot n \cdot 2^n
  \quad \Rightarrow \quad T(n) \in \BigTheta{n \cdot 2^n} \qquad \square
\end{equation}

\subsection{Tabla comparativa de complejidades}

\begin{table}[ht]
\centering
\caption{Complejidades de los componentes de K-GeoMIP}
\label{tab:complejidades}
\begin{tabular}{@{}l l l l@{}}
\toprule
\textbf{Método} & \textbf{Tiempo} & \textbf{Espacio} & \textbf{Notas} \\
\midrule
\texttt{System.bipartir}        & $\BigO{n \cdot 2^d}$ & $\BigO{n \cdot 2^d}$ & $k=2$ fijo \\
\texttt{System.k\_partir}       & $\BigO{n \cdot 2^d}$ & $\BigO{n \cdot 2^d}$ & $k \geq 2$ \\
\texttt{stirling(n,k)}          & $\BigTheta{n \cdot k}$ & $\BigO{k}$ & DP 2 filas \\
\texttt{particionar\_conjunto}  & $\BigTheta{n \cdot \Stir{n}{k}}$ & $\BigO{n}$ & RGS lazy \\
\texttt{KGeometricSIA.aplicar}  & $\BigTheta{n \cdot 2^n}$ & $\BigO{n \cdot 2^n}$ & greedy $C=25$ \\
Búsqueda exhaustiva             & $\BigO{\Stir{n}{k}^2 \cdot n \cdot 2^n}$ &
  $\BigO{n \cdot 2^n}$ & inviable $n \geq 8$ \\
\bottomrule
\end{tabular}
\end{table}

\subsection{Speedup respecto a búsqueda exhaustiva}

\begin{equation}
  \mathrm{Speedup}(n, k) = \frac{\Stir{n}{k}^2}{C}
  \quad \xrightarrow{n=10,\; k=3} \quad
  \frac{9\,330^2}{25} \approx \mathbf{3{,}47 \times 10^6}
  \label{eq:speedup}
\end{equation}

% ════════════════════════════════════════════════════════════
%  SECCIÓN 6 — DETALLES DE IMPLEMENTACIÓN
% ════════════════════════════════════════════════════════════
\section{Detalles de Implementación}
\label{sec:implementacion}

\subsection{Métodos principales de \texttt{KGeometricSIA}}

\begin{table}[ht]
\centering
\caption{API pública y privada de \texttt{KGeometricSIA}}
\label{tab:api-kgeomip}
\begin{tabular}{@{}l p{5.5cm} l@{}}
\toprule
\textbf{Método} & \textbf{Descripción} & \textbf{Retorno} \\
\midrule
\texttt{aplicar\_estrategia(cond,alc,mec,tpm,k)} &
  Orquesta el flujo completo & \texttt{Solution} \\
\texttt{\_generar\_candidatos\_k(k)} &
  Genera hasta $C=25$ candidatos & \texttt{list[Partición]} \\
\texttt{\_greedy\_multiway(c, n, k, desc)} &
  Algoritmo LPT & $k$ grupos \\
\texttt{\_grupos\_a\_particion(grupos, A, M, k)} &
  Convierte índices a NDArray & \texttt{list[tuple]} \\
\texttt{\_particion\_valida(p)} &
  Verifica partes no vacías & \texttt{bool} \\
\texttt{\_particion\_a\_clave(p)} &
  Clave hashable canónica & \texttt{frozenset} \\
\texttt{\_dist\_particion(p)} &
  Dist.\ marginal con caché (Opt 6) & \texttt{NDArray} \\
\texttt{\_particion\_a\_grupos(p, A)} &
  Asignación a índices locales & \texttt{list[list]} \\
\texttt{\_refinar\_local(p, emd, dist)} &
  Búsqueda local (Opt 7) & \texttt{tuple} \\
\bottomrule
\end{tabular}
\end{table}

\subsection{Optimizaciones NumPy}

\subsubsection{Optimización 1 --- Dict a NDArray 2D}

\begin{lstlisting}[style=python,
  caption={Matriz de costos: dict vs NDArray},label=lst:opt1]
# ANTES -- dict con claves tupla, O(n) hash por acceso
tabla_transiciones = {}  # dict[tuple, list]

# DESPUES -- ndarray 2D, O(1) por indice entero
self._trans_matrix = np.zeros((1 << n_mec, n_fut), dtype=np.float32)
self._trans_valid  = np.zeros(1 << n_mec, dtype=bool)

@staticmethod
def _estado_a_idx(estado) -> int:
    # Little-endian: estado=[1,0,0] -> idx=1
    return int("".join(map(str, reversed(list(estado)))), 2)
\end{lstlisting}

\subsubsection{Optimización 2 --- Vectorización de acumulación BFS}

\begin{lstlisting}[style=python,
  caption={Acumulación vectorizada en el BFS},label=lst:opt2]
# ANTES -- dos bucles Python O(n_fut) cada uno
for n in ncubos:
    tabla[key][n] += tabla[tmp_key][n]
tmp = [factor * x for x in tabla[key]]

# DESPUES -- operaciones SIMD via NumPy, O(n_fut) en C
self._trans_matrix[fin_idx] += self._trans_matrix[tmp_idx]
self._trans_matrix[fin_idx] *= factor
\end{lstlisting}

\subsubsection{Optimización 3 --- Vectorización de identificación de particiones}

\begin{lstlisting}[style=python,
  caption={Identificación vectorizada de particiones óptimas},label=lst:opt3]
# ANTES -- bucles Python
presentes = [idx for idx, i in enumerate(estado)
             if i == self.caminos[0][0][idx]]

# DESPUES -- NumPy
estado_arr = np.array(estado)
ini_arr    = np.array(self.caminos[0][0])
presentes  = np.where(estado_arr == ini_arr)[0].tolist()
costos_min = np.minimum(actual, complementario)
futuros    = np.where(actual <= complementario)[0].tolist()
costo      = float(costos_min.sum())
\end{lstlisting}

El speedup estimado por las tres optimizaciones combinadas para $n=10$:
\begin{equation}
  S \approx \frac{t_{\mathrm{CPython}}}{t_{\mathrm{NumPy\;SIMD}}}
  \approx \frac{n \cdot 2^n \cdot 200\,\mathrm{ns}}{n \cdot 2^n \cdot 5\,\mathrm{ns}}
  \approx 40\times
\end{equation}

\subsection{Análisis de cuellos de botella y optimizaciones v1.1}
\label{sec:cuellos-botella}

Tras un perfilado del flujo completo de \texttt{KGeometricSIA} sobre los sistemas
$N10A$ y $N15A$ se identificaron cuatro cuellos de botella, ordenados por impacto.
La Tabla~\ref{tab:bottlenecks} los resume junto con la optimización aplicada a cada uno.

\begin{table}[ht]
\centering
\caption{Cuellos de botella detectados y optimizaciones v1.1}
\label{tab:bottlenecks}
\begin{tabular}{@{}l p{4.6cm} l l@{}}
\toprule
\textbf{\#} & \textbf{Cuello de botella} & \textbf{Fase} & \textbf{Optimización} \\
\midrule
B1 & \texttt{\_estado\_a\_idx} basado en \texttt{str} ($\Theta(n\,2^n)$ llamadas) &
  BFS (2--3) & Opt 4 \\
B2 & \texttt{calcular\_costo}: listas Python por estado &
  BFS (2--3) & Opt 5 \\
B3 & Re-marginalización completa por candidato en Fase 5 &
  Evaluación (5) & Opt 6 \\
B4 & Candidatos greedy sin refinamiento (calidad de $\delta_k$) &
  Selección (6) & Opt 7 \\
\bottomrule
\end{tabular}
\end{table}

\subsubsection{Optimización 4 --- \texttt{\_estado\_a\_idx} con aritmética de bits}

La conversión estado~$\to$~índice entero (convenio little-endian) se invoca
$\BigTheta{n \cdot 2^n}$ veces durante la expansión BFS. La implementación previa
(Listado~\ref{lst:opt1}) construía y parseaba una cadena por llamada, generando
basura intermedia (\texttt{list}~$\to$~\texttt{map}~$\to$~\texttt{join}~$\to$~\texttt{int})
en el bucle más caliente del algoritmo. Se reemplazó por aritmética de bits pura,
con resultado idéntico bit a bit (el bit en la posición $i$ conserva su peso $2^i$).

\begin{lstlisting}[style=python,
  caption={Opt 4: estado a indice sin construir strings},label=lst:opt4]
# ANTES (v1.0) -- construye y parsea un string por cada estado
return int("".join(map(str, reversed(list(estado)))), 2)

# DESPUES (v1.1) -- OR de bits, sin asignaciones intermedias
idx = 0
for posicion, bit in enumerate(estado):
    if bit:
        idx |= 1 << posicion   # activa el bit i con peso 2^i (little-endian)
return idx
\end{lstlisting}

\subsubsection{Optimización 5 --- Vectorización de \texttt{calcular\_costo} con \texttt{\_flat\_matrix}}

Por cada uno de los $2^n$ estados, \texttt{calcular\_costo} reconstruía dos listas
Python \texttt{[flat[idx] for flat in \_flat\_data]} para formar el vector de diferencias
$|X_{\mathrm{ini}} - X_{\mathrm{fin}}|$. Se precomputa una sola vez la matriz
$\texttt{\_flat\_matrix} \in \mathbb{R}^{n_{\mathrm{fut}} \times 2^{n_{\mathrm{mec}}}}$
(apilando los $n$-cubos aplanados) y se extraen las columnas por \emph{fancy indexing},
eliminando $2 \cdot n_{\mathrm{fut}}$ accesos indexados en Python por estado.

\begin{lstlisting}[style=python,
  caption={Opt 5: diferencias por columnas vectorizadas},label=lst:opt5]
# Precomputo (una vez, junto a _flat_data):
self._flat_matrix = np.stack(self._flat_data)   # (n_fut, 2^n_mec)

# ANTES (v1.0) -- dos comprensiones Python por estado
diffs = np.abs(
    np.array([flat[ini] for flat in self._flat_data])
    - np.array([flat[fin] for flat in self._flat_data]))

# DESPUES (v1.1) -- extraccion de columnas en C, sin listas temporales
diffs = np.abs(self._flat_matrix[:, ini] - self._flat_matrix[:, fin])
\end{lstlisting}

Las optimizaciones 4 y 5 \textbf{no alteran el resultado}: se verificó la identidad
bit a bit de \texttt{\_estado\_a\_idx} para todas las longitudes $n \in [1,12]$ y
$200$ estados aleatorios por longitud. Su efecto es exclusivamente temporal sobre la
fase dominante $\BigTheta{n \cdot 2^n}$, medido en la Tabla~\ref{tab:opt-tiempos}.

\begin{table}[ht]
\centering
\caption{Efecto temporal de Opt 4--6 sobre la fase BFS (resultado idéntico)}
\label{tab:opt-tiempos}
\begin{tabular}{@{}l l r r r@{}}
\toprule
\textbf{Sistema} & \textbf{Fase} & \textbf{v1.0 (s)} & \textbf{v1.1 (s)} & \textbf{Mejora} \\
\midrule
$N15A$, $k=2$ & BFS completo & $5{,}55$ & $4{,}38$ & $-21\%$ \\
$N15A$, $k=3$ & BFS + Fase 5 & $5{,}34$ & $4{,}29$ & $-20\%$ \\
$N10A$, $k=2$ & BFS completo & $0{,}124$ & $0{,}112$ & $-10\%$ \\
\bottomrule
\end{tabular}
\end{table}

\subsubsection{Optimización 6 --- Caché de marginales por (cubo, mecanismo)}

En la Fase 5 cada candidato ejecutaba
\texttt{k\_partir(P).distribucion\_marginal()}, reconstruyendo un \texttt{System}
completo y re-marginalizando \emph{todos} los $n$-cubos. La distribución marginal de una
k-partición, sin embargo, se factoriza por $n$-cubo: el valor del cubo $i$ depende
\textbf{únicamente} de su índice y del conjunto-mecanismo $M_j$ de la parte que lo
contiene en su alcance, y \emph{no} del resto de la partición:
\begin{equation}
  \mathrm{dist}[i] \;=\; 1 - P\!\Big(\texttt{marginalizar}\big(c_i,\; c_i.\mathrm{dims} \setminus M_j\big)\Big),
  \qquad c_i.\mathrm{indice} \in A_j .
\end{equation}
Por ello se cachea el valor marginal por la clave $(c_i.\mathrm{indice},\, \mathrm{frozenset}(M_j))$,
evitando repetir el costoso \texttt{np.mean} sobre $2^d$ celdas cuando el mismo par
$(\text{cubo},\text{mecanismo})$ reaparece entre candidatos o entre movimientos de la
búsqueda local. Se verificó que \texttt{\_dist\_particion} produce un resultado
\textbf{idéntico} ($\max|\Delta| = 0{,}0$) al de \texttt{k\_partir().distribucion\_marginal()}
sobre todos los candidatos de $N10A$ con $k \in \{3,4,5\}$.

\subsection{Sistema de caché de subsistema (v0.9)}

Clave de caché: $\mathrm{clave} = (\mathrm{condicion}, \mathrm{alcance}, \mathrm{mecanismo})$.

\begin{center}
\begin{tabular}{@{}l l@{}}
\toprule
\textbf{Escenario} & \textbf{Costo total Fases 1--3} \\
\midrule
Antes de v0.9 ($k \in \{2,3,4,5\}$, mismo subsistema) & $4 \cdot \BigTheta{n \cdot 2^n}$ \\
Después de v0.9 (1 miss + 3 hits) & $1 \cdot \BigTheta{n \cdot 2^n}$ \\
Ahorro absoluto & $3 \cdot \BigTheta{n \cdot 2^n}$ (\textbf{75\%}) \\
\bottomrule
\end{tabular}
\end{center}

\subsection{Suite de tests de validación}

\begin{table}[ht]
\centering
\caption{Suite de 13 tests automatizados}
\label{tab:tests}
\begin{tabular}{@{}l l l l@{}}
\toprule
\textbf{Grupo} & \textbf{Tests} & \textbf{Postcondición} & \textbf{Dataset} \\
\midrule
A -- Corrección formal  & 2 &
  $|\delta_2^{\mathrm{KGeo}} - \delta_2^{\mathrm{Geo}}| < 10^{-9}$;
  \texttt{ValueError} para $k \notin [2,5]$ & N4A \\
B -- Validez solución   & 6 &
  pérdida $\geq 0$, distribuciones $\neq$ None, tiempo $\geq 0$ & N6A \\
C -- Rendimiento        & 5 &
  speedup greedy, sin regresión, \texttt{stirling} correcto & N6A \\
\midrule
\textbf{Total} & \textbf{13/13} & \textbf{passed en 7{,}13 s} & Python 3.12.5 \\
\bottomrule
\end{tabular}
\end{table}

\subsection{Dependencias externas}

\begin{center}
\begin{tabular}{@{}l l l@{}}
\toprule
\textbf{Biblioteca} & \textbf{Versión} & \textbf{Uso} \\
\midrule
NumPy   & $\geq$ 2.0  & Arrays n-dimensionales, operaciones SIMD \\
SciPy   & $\geq$ 1.10 & EMD (Wasserstein), cálculo científico \\
pandas  & $\geq$ 2.0  & Lectura/escritura de Excel (.xlsx) \\
pytest  & $\geq$ 9.0  & Framework de testing automatizado \\
colorama & $\geq$ 0.4 & Salida en consola con colores (Windows) \\
\bottomrule
\end{tabular}
\end{center}

% ════════════════════════════════════════════════════════════
%  SECCIÓN 7 — RESULTADOS EXPERIMENTALES
% ════════════════════════════════════════════════════════════
\section{Resultados Experimentales}
\label{sec:resultados}

\subsection{Entorno experimental}

\begin{center}
\begin{tabular}{@{}l l@{}}
\toprule
\textbf{Parámetro} & \textbf{Valor} \\
\midrule
Procesador & Intel Core i5-12500H (12.ª Gen) \\
RAM        & 32\,GB \\
S.O.       & Windows 11 \\
Python     & 3.12.5 \\
NumPy      & 2.4.4 \\
Referencia & \texttt{GeometricSIA} ($k=2$, búsqueda exacta, $\BigTheta{n \cdot 2^n}$) \\
Evaluado   & \texttt{KGeometricSIA} ($k \in \{3,4,5\}$, greedy, $\BigTheta{n \cdot 2^n}$) \\
Conjunto oficial & \texttt{DatosPruebas2026\_1.xlsx} (hojas 10A, 15B, 20A, 22A, 25A) \\
Microbenchmark & Datasets incrementales N3A--N10A (escalabilidad temporal) \\
Estado inicial & \texttt{"1"} seguido de $n-1$ ceros \\
\bottomrule
\end{tabular}
\end{center}

Los resultados de calidad ($\delta_k$) provienen del \textbf{conjunto oficial de pruebas}
\texttt{DatosPruebas2026\_1.xlsx}, que evalúa $\approx 50$ subsistemas por sistema. Dicho
archivo registró \emph{partición} y \emph{pérdida} pero \emph{no} tiempos de ejecución;
por ello los tiempos de las Tablas~\ref{tab:tiempos} y~\ref{tab:ratio} provienen de un
microbenchmark controlado sobre los datasets incrementales N3A--N10A (un subsistema completo
por tamaño), medido en la misma plataforma.

\subsection{Tabla 1 --- Tiempos de ejecución (segundos)}

\begin{table}[ht]
\centering
\caption{Tiempo total de \texttt{aplicar\_estrategia} en segundos}
\label{tab:tiempos}
\begin{tabular}{@{}r r r r r@{}}
\toprule
$n$ & $k=2$ (GeoMIP) & $k=3$ & $k=4$ & $k=5$ \\
\midrule
3  & 0{,}0117 & 0{,}0086 & ---    & ---    \\
4  & 0{,}0100 & 0{,}0119 & 0{,}0128 & ---    \\
5  & 0{,}0173 & 0{,}0163 & 0{,}0164 & 0{,}0160 \\
6  & 0{,}0391 & 0{,}0240 & 0{,}0197 & 0{,}0247 \\
8  & 0{,}0537 & 0{,}0488 & 0{,}0372 & 0{,}0726 \\
10 & 0{,}1682 & 0{,}1074 & 0{,}1313 & 0{,}1216 \\
\bottomrule
\end{tabular}
\end{table}

\subsection{Tabla 2 --- Ratio temporal $t_k / t_{\mathrm{geo}}$}

\begin{table}[ht]
\centering
\caption{Ratio KGeometricSIA$(k)$ vs.\ GeometricSIA$(k=2)$}
\label{tab:ratio}
\begin{tabular}{@{}r r r r@{}}
\toprule
$n$ & $k=3\,/\,k=2$ & $k=4\,/\,k=2$ & $k=5\,/\,k=2$ \\
\midrule
3  & $0{,}74\times$ & --- & --- \\
4  & $1{,}20\times$ & $1{,}28\times$ & --- \\
5  & $0{,}94\times$ & $0{,}95\times$ & $0{,}92\times$ \\
6  & $0{,}61\times$ & $0{,}50\times$ & $0{,}63\times$ \\
8  & $0{,}91\times$ & $0{,}69\times$ & $1{,}35\times$ \\
10 & $0{,}64\times$ & $0{,}78\times$ & $0{,}72\times$ \\
\bottomrule
\end{tabular}
\end{table}

Ratio máximo: $1{,}35\times$ ($n=8$, $k=5$) $\ll 3{,}0\times$ (umbral del test C2).

\subsection{Tabla 3 --- Calidad de la solución ($\delta_k$ EMD) sobre el conjunto oficial}

La Tabla~\ref{tab:calidad} reporta la pérdida EMD \textbf{promedio} de \texttt{KGeometricSIA}
(columna \emph{Geometric} del archivo oficial) sobre los $\approx 50$ subsistemas de cada hoja
de \texttt{DatosPruebas2026\_1.xlsx}. Cada celda es la media de $\delta_k$ sobre todos los
subsistemas evaluados de ese sistema.

\begin{table}[ht]
\centering
\caption{Pérdida EMD media $\overline{\delta_k}$ por sistema (conjunto oficial)}
\label{tab:calidad}
\begin{tabular}{@{}r r r r r r@{}}
\toprule
$n$ & \#subsist. & $\overline{\delta_2}$ & $\overline{\delta_3}$ & $\overline{\delta_4}$ & $\overline{\delta_5}$ \\
\midrule
10 (10A) & 49 & 0{,}0766 & 1{,}3583 & 1{,}3906 & 1{,}4449 \\
15 (15B) & 50 & 0{,}0002 & 0{,}0166 & 0{,}0198 & 0{,}0213 \\
20 (20A) & 50 & 0{,}0702 & 2{,}6482 & 2{,}6631 & --- \\
22 (22A) & 50 & \multicolumn{4}{c}{sin completar (presupuesto temporal excedido)} \\
25 (25A) & 50 & \multicolumn{4}{c}{sin completar (presupuesto temporal excedido)} \\
\bottomrule
\end{tabular}
\end{table}

\noindent\textbf{Observaciones:}
\begin{itemize}
  \item La pérdida media crece monótonamente con $k$
    ($\overline{\delta_2} \leq \overline{\delta_3} \leq \overline{\delta_4} \leq \overline{\delta_5}$)
    en todos los sistemas, coherente con que más particiones implican más pérdida de información.
  \item El sistema $N15B$ presenta una pérdida casi nula en $k=2$
    ($\overline{\delta_2} = 0{,}0002$): el subsistema es prácticamente reducible por bipartición.
  \item Los sistemas $N22A$ y $N25A$ no se completaron por el costo $\BigTheta{n \cdot 2^n}$ de la
    fase BFS ($2^{22}$ y $2^{25}$ estados); su evaluación es objetivo de las optimizaciones de la
    sección~\ref{sec:cuellos-botella} y del trabajo futuro (Tabla~\ref{tab:mejoras}).
\end{itemize}

\subsubsection*{Comparativa con QNodes ($\overline{\delta_k}$)}

El archivo oficial incluye además los resultados de la estrategia \texttt{QNodes} en los
sistemas donde se completó. La Tabla~\ref{tab:calidad-qnodes} contrasta ambas medias.

\begin{table}[ht]
\centering
\caption{$\overline{\delta_k}$: \texttt{QNodes} vs.\ \texttt{Geometric} (K-GeoMIP)}
\label{tab:calidad-qnodes}
\begin{tabular}{@{}l l r r r@{}}
\toprule
\textbf{Sistema} & \textbf{Estrategia} & $\overline{\delta_2}$ & $\overline{\delta_3}$ & $\overline{\delta_4}$ \\
\midrule
\multirow{2}{*}{$N10A$ (49)} & QNodes    & 0{,}0982 & \textbf{0{,}2778} & \textbf{0{,}4597} \\
                             & Geometric & \textbf{0{,}0766} & 1{,}3583 & 1{,}3906 \\
\midrule
\multirow{2}{*}{$N15B$ (50)} & QNodes    & 0{,}0002 & \textbf{0{,}0016} & --- \\
                             & Geometric & 0{,}0002 & 0{,}0166 & 0{,}0198 \\
\bottomrule
\end{tabular}
\end{table}

\noindent En $k=2$ la búsqueda exacta de \texttt{Geometric} iguala o mejora a \texttt{QNodes}
(p.\,ej.\ $N10A$: $0{,}0766$ vs.\ $0{,}0982$). En $k \geq 3$, en cambio, \texttt{QNodes} obtiene
particiones de menor pérdida media: la heurística greedy de K-GeoMIP explora solo $C=25$
candidatos de un espacio de $S(n,k)^2$ particiones, mientras que QNodes aprovecha la estructura
de nodos de corte. Esta brecha motivó directamente la \textbf{Optimización~7} (búsqueda local,
sección~\ref{sec:busqueda-local}), que reduce $\delta_k$ hasta un $-14{,}7\%$ sin regresión;
cerrar por completo la brecha es objeto de las mejoras propuestas (Beam Search, Branch \& Bound)
en la Tabla~\ref{tab:mejoras}.

\subsection{Tabla 4 --- Reducción del espacio de búsqueda}

\begin{table}[ht]
\centering
\caption{Speedup de la fase de búsqueda: $\Stir{n}{k}^2 / C$}
\label{tab:speedup}
\begin{tabular}{@{}r r r r r r r@{}}
\toprule
$n$ & $\Stir{n}{3}$ & Sp $k=3$ & $\Stir{n}{4}$ & Sp $k=4$ & $\Stir{n}{5}$ & Sp $k=5$ \\
\midrule
5  & 25    & $8{,}3\times$      & 10    & $3{,}3\times$    & 1     & $0{,}3\times$ \\
6  & 90    & $\mathbf{30\times}$ & 65    & $\mathbf{21\times}$ & 15    & $\mathbf{5\times}$ \\
8  & 966   & $\mathbf{322\times}$ & 1701  & $\mathbf{567\times}$ & 1050  & $\mathbf{350\times}$ \\
10 & 9330  & $\mathbf{3110\times}$ & 34105 & $\mathbf{11368\times}$ & 42525 & $\mathbf{14175\times}$ \\
\bottomrule
\end{tabular}
\end{table}

\subsection{Gráfica de escalabilidad temporal}

\begin{figure}[ht]
\centering
\begin{tikzpicture}
\begin{axis}[
  width=12cm, height=7cm,
  xlabel={Número de variables $n$},
  ylabel={Tiempo de ejecución (s)},
  xmin=2, xmax=11,
  ymin=0, ymax=0.2,
  legend pos=north west,
  grid=major,
  grid style={dashed, gray!30},
  xtick={3,4,5,6,8,10},
  title={Escalabilidad temporal de K-GeoMIP vs.\ GeoMIP},
]
\addplot[color=KGeoBlue, mark=*, thick]
  coordinates {(3,0.0117)(4,0.0100)(5,0.0173)(6,0.0391)(8,0.0537)(10,0.1682)};
\addlegendentry{$k=2$ (GeoMIP)}
\addplot[color=KGeoGreen, mark=square*, thick]
  coordinates {(3,0.0086)(4,0.0119)(5,0.0163)(6,0.0240)(8,0.0488)(10,0.1074)};
\addlegendentry{$k=3$ (K-GeoMIP)}
\addplot[color=KGeoOrange, mark=triangle*, thick]
  coordinates {(4,0.0128)(5,0.0164)(6,0.0197)(8,0.0372)(10,0.1313)};
\addlegendentry{$k=4$ (K-GeoMIP)}
\addplot[color=violet, mark=diamond*, thick]
  coordinates {(5,0.0160)(6,0.0247)(8,0.0726)(10,0.1216)};
\addlegendentry{$k=5$ (K-GeoMIP)}
\end{axis}
\end{tikzpicture}
\caption{Tiempos de ejecución para $n \in \{3,4,5,6,8,10\}$.
  Todos los métodos exhiben crecimiento $\BigTheta{n \cdot 2^n}$.}
\label{fig:tiempos}
\end{figure}

% ════════════════════════════════════════════════════════════
%  SECCIÓN 8 — LIMITACIONES Y TRABAJO FUTURO
% ════════════════════════════════════════════════════════════
\section{Limitaciones y Trabajo Futuro}
\label{sec:limitaciones}

\subsection{Limitaciones de la heurística greedy}

La heurística \textbf{no} garantiza encontrar la partición óptima global. Tres razones:
\begin{enumerate}
  \item \textbf{Cobertura limitada:} Para $n=10$, $k=3$: se evalúan 25 de
    $\approx 87$M particiones posibles ($0{,}027\%$).
  \item \textbf{Vector $\mathbf{c}$ como proxy del EMD:} $c_i$ captura diferencias
    marginales de un nodo; el EMD real captura correlaciones conjuntas entre nodos.
  \item \textbf{Garantía LPT $\neq$ garantía EMD:}
    La cota (\ref{eq:lpt-garantia}) es sobre makespan, no sobre $\delta_k$.
\end{enumerate}

\subsection{Refinamiento por búsqueda local (Optimización 7, v1.1)}
\label{sec:busqueda-local}

Para mitigar la limitación~1 (cobertura limitada del espacio de búsqueda) se añadió
una fase de \textbf{búsqueda local de primera-mejora} (\emph{hill climbing}) que pule la
mejor k-partición hallada por el greedy. La solución se representa como una asignación de
las $n$ variables a $k$ grupos; en cada pase se intenta mover cada variable a otro grupo,
se reconstruye y valida la k-partición, y se evalúa su $\delta_k$. Se acepta el primer
movimiento que reduce estrictamente $\delta_k$ y se reinicia el pase desde la nueva solución.

\begin{lstlisting}[style=pseudocode,
  caption={Refinamiento por busqueda local (first-improvement)},label=lst:opt7]
refinar_local(particion, emd, dist):
    si n > REFINAMIENTO_N_MAX: retornar (particion, emd, dist)   # gate de tamano
    grupos <- asignacion_local(particion)
    mientras hubo_mejora y evals < MAX_REFINAMIENTO_EVALS:
        para cada grupo gi, cada variable v en gi, cada destino gj != gi:
            mover v de gi a gj
            cand <- grupos_a_particion(grupos)
            si cand es valida:
                d <- dist_particion(cand);  e <- emd_efecto(d, dist_orig)   # Opt 6
                si e < mejor_emd - eps:  aceptar y reiniciar pase
                si no: revertir movimiento
            si no: revertir movimiento
    retornar mejor (particion, emd, dist)
\end{lstlisting}

\paragraph{Garantía de no-regresión (mejora monótona).}
La partición devuelta cumple $\delta_k^{\,\mathrm{v1.1}} \leq \delta_k^{\,\mathrm{v1.0}}$
\emph{siempre}. La clave es que el caché inter-$k$ (warm-start) guarda la mejor partición
\textbf{antes} del refinamiento; por tanto la generación de candidatos de cada $k$ es
idéntica a la de v1.0 y el pool siempre contiene la solución base previa. Como solo se
aceptan movimientos de mejora estricta, el refinamiento únicamente puede igualar o reducir
la pérdida, nunca aumentarla.

\paragraph{Acotación de costo.}
El refinamiento se \textbf{desactiva} si $n > \texttt{REFINAMIENTO\_N\_MAX} = 16$, de modo
que los sistemas grandes solo reciben el \emph{speedup} del BFS (Opt 4--5) y nunca trabajo
adicional. Para $n \leq 16$ el número de evaluaciones está acotado por
$\texttt{MAX\_REFINAMIENTO\_EVALS} = 40$, y cada evaluación reutiliza el caché de marginales
(Opt 6): mover una sola variable solo recomputa los marginales de los grupos afectados.
El costo extra $\BigO{40 \cdot n \cdot 2^d}$ queda dominado por la fase BFS $\BigTheta{n \cdot 2^n}$.

\begin{table}[ht]
\centering
\caption{Reducción de pérdida $\delta_k$ por el refinamiento local (sin regresión)}
\label{tab:opt7-perdida}
\begin{tabular}{@{}l c r r r@{}}
\toprule
\textbf{Sistema} & \textbf{$k$} & \textbf{$\delta_k$ v1.0} & \textbf{$\delta_k$ v1.1} & \textbf{Mejora} \\
\midrule
$N10A$ & 3 & $4{,}34375$ & $3{,}70703$ & $-14{,}7\%$ \\
$N10A$ & 4 & $4{,}37891$ & $4{,}10547$ & $-6{,}2\%$ \\
$N10A$ & 5 & $4{,}41797$ & $4{,}41797$ & $0\%$ \\
$N15A$ & 3 & $0{,}341422$ & $0{,}322286$ & $-5{,}6\%$ \\
$N15A$ & 4 & $0{,}358117$ & $0{,}343485$ & $-4{,}1\%$ \\
$N15A$ & 5 & $0{,}369636$ & $0{,}366972$ & $-0{,}7\%$ \\
\bottomrule
\end{tabular}
\end{table}

\subsection{Comparativa de estrategias de búsqueda}

\begin{table}[ht]
\centering
\caption{Alternativas de búsqueda para k-MIP}
\label{tab:alternativas}
\begin{tabular}{@{}l l l l@{}}
\toprule
\textbf{Estrategia} & \textbf{Garantía} & \textbf{Complejidad} & \textbf{$n$ viable} \\
\midrule
Exhaustiva        & Exacta (100\%) & $\BigO{\Stir{n}{k}^2 \cdot 2^n}$ & $n \leq 6$ \\
Greedy $C=25$ (actual) & Ninguna formal & $\BigO{n \log n}$ & $n \leq 25$ \\
Branch and Bound  & Exacta (100\%) & $\BigO{k^d}$ (con poda) & $n \leq 14$ práctico \\
Beam Search $B=50$ & Cercana & $\BigO{B \cdot n \cdot k}$ & $n \leq 20$ \\
\bottomrule
\end{tabular}
\end{table}

La cota inferior del B\&B es:
\begin{equation}
  \mathrm{lb}(\mathrm{estado}) =
  \mathrm{costo\_actual} + \frac{1}{k} \sum_{j \in \mathrm{pendientes}} c_j
\end{equation}

\subsection{Mejoras futuras propuestas}

\begin{table}[ht]
\centering
\caption{Mejoras propuestas para v2.0}
\label{tab:mejoras}
\begin{tabular}{@{}l p{4cm} l l l@{}}
\toprule
\textbf{\#} & \textbf{Mejora} & \textbf{Calidad} & \textbf{Velocidad} & \textbf{Esfuerzo} \\
\midrule
17.1 & Búsqueda local first-improvement & Medio & Negativo leve & --- \textbf{(hecho v1.1)} \\
17.2 & Branch \& Bound ($n \leq 12$) & Alto & Negativo & Alto \\
17.3 & Beam Search ($W$ configurable) & Medio-Alto & Neutro & Medio \\
17.4 & EMD exacta en Fase A & Medio & Neutro & Bajo \\
17.5 & Caché persistente en disco & Ninguno & Muy alto & Bajo \\
17.6 & Particiones no balanceadas & Medio & Neutro & Medio \\
17.7 & Paralelismo Fase 5 ($\approx3\times$) & Ninguno & Alto & Bajo \\
\bottomrule
\end{tabular}
\end{table}

La mejora 17.1 (búsqueda local) ya está implementada en v1.1; véase la
sección~\ref{sec:busqueda-local}. \textbf{Prioridad recomendada para v2.0:}
$17.5 \to 17.7 \to 17.4 \to 17.3$ (máximo beneficio con mínimo esfuerzo).

% ════════════════════════════════════════════════════════════
%  SECCIÓN 9 — USO DE IA GENERATIVA
% ════════════════════════════════════════════════════════════
\section{Uso de Inteligencia Artificial Generativa}

Durante el desarrollo del proyecto K-GeoMIP se emplearon herramientas de IA generativa
de forma transparente y complementaria al trabajo propio:

\begin{itemize}
  \item \textbf{Claude Code (Anthropic):} Asistencia en la estructuración del pseudocódigo
    del algoritmo RGS y en la generación de la suite de tests automatizados.
    Prompts típicos: ``\textit{¿Cómo implementar un generador RGS en Python?}'',
    ``\textit{Genera tests parametrizados para validar equivalencia con GeometricSIA}''.

  \item \textbf{GitHub Copilot:} Autocompletado de código repetitivo en operaciones
    NumPy y generación de docstrings.

  \item \textbf{Influencia en decisiones de diseño:} La idea de usar \texttt{frozenset}
    como clave de caché para detectar duplicados surgió de una sugerencia de IA,
    validada experimentalmente.
\end{itemize}

Todas las decisiones algorítmicas críticas (selección de LPT, demostración de equivalencia
\texttt{k\_partir} $\equiv$ \texttt{bipartir}, análisis $\BigTheta{n \cdot 2^n}$) fueron
verificadas independientemente por el autor.

% ════════════════════════════════════════════════════════════
%  APÉNDICES
% ════════════════════════════════════════════════════════════
\appendix

\section{Demostraciones}
\label{apx:demostraciones}

\subsection{Corrección del algoritmo RGS}

\begin{lemma}[Completitud sin duplicados]
El algoritmo RGS genera exactamente todas las k-particiones distintas de
$\{0, \ldots, n-1\}$, sin repeticiones.
\end{lemma}

\begin{proof}
\textbf{Completitud:} Toda k-partición tiene exactamente una RGS canónica con $a[0] = 0$.
El algoritmo recorre todas las RGS válidas por DFS $\Rightarrow$ completitud.

\textbf{Sin duplicados:} Dos k-particiones son iguales ssi tienen la misma RGS canónica.
El DFS genera cada RGS exactamente una vez $\Rightarrow$ sin duplicados. $\square$
\end{proof}

\subsection{Invariante de factorización de \texttt{k\_partir}}

\begin{lemma}[Propiedad de factorización]
Sea $P = \texttt{k\_partir}(\mathrm{particion})$. Entonces:
\begin{equation}
  \mathrm{dist\_marginal}(P) = \bigotimes_{i=1}^{k} p_{S_i}(A_i, M_i)
\end{equation}
\end{lemma}

\begin{proof}
Cada n-cubo $c$ en $A_i$ queda marginalizado sobre $\mathrm{dims} \setminus M_i$,
equivalente a condicionar sobre solo las variables $M_i$.
Por independencia condicional entre las $k$ partes, la distribución conjunta es el
producto tensorial. $\square$
\end{proof}

\section{Resultados de Tests}
\label{apx:tests}

\begin{lstlisting}[style=pseudocode,
  caption={Salida de la suite de 13 tests},label=lst:tests]
platform win32 -- Python 3.12.5, pytest-9.0.3
tests/test_kgeomip.py::test_k2_equivalencia_geometrica PASSED
tests/test_kgeomip.py::test_k_invalido_lanza_valueerror[k=0] PASSED
tests/test_kgeomip.py::test_k_invalido_lanza_valueerror[k=1] PASSED
tests/test_kgeomip.py::test_k_invalido_lanza_valueerror[k=6] PASSED
tests/test_kgeomip.py::test_solucion_atributos_validos[k=3] PASSED
tests/test_kgeomip.py::test_solucion_atributos_validos[k=4] PASSED
tests/test_kgeomip.py::test_solucion_atributos_validos[k=5] PASSED
tests/test_kgeomip.py::test_distribucion_es_ndarray[k=3] PASSED
tests/test_kgeomip.py::test_distribucion_es_ndarray[k=4] PASSED
tests/test_kgeomip.py::test_distribucion_es_ndarray[k=5] PASSED
tests/test_kgeomip.py::test_aislamiento_fase_busqueda PASSED
tests/test_kgeomip.py::test_sin_regresion_rendimiento PASSED
tests/test_kgeomip.py::test_stirling_correctness[k=3] PASSED
========== 13 passed, 32 warnings in 7.13s ==========
\end{lstlisting}

\section{Registro de Cambios}
\label{apx:changelog}

\begin{longtable}{@{}l l p{8cm}@{}}
\toprule
\textbf{Versión} & \textbf{Fecha} & \textbf{Cambio} \\
\midrule
\endhead
v1.1 & 2026-06-15 &
  Revisión de cuellos de botella (sección~\ref{sec:cuellos-botella}):
  Opt~4 (\texttt{\_estado\_a\_idx} con aritmética de bits) y
  Opt~5 (\texttt{calcular\_costo} vectorizado con \texttt{\_flat\_matrix}) reducen
  $-20\%$ la fase BFS sin alterar el resultado;
  Opt~6 (caché de marginales por (cubo, mecanismo)) y
  Opt~7 (búsqueda local first-improvement, sección~\ref{sec:busqueda-local})
  reducen $\delta_k$ hasta $-14{,}7\%$ con garantía de no-regresión. \\[4pt]
v1.0 & 2026-06-14 &
  Optimización por chunks en \texttt{Manager.generar\_red}:
  RAM pico $O(2^n \times n) \to O(65{,}536 \times n) \approx 4\,\text{MB}$ constante. \\[4pt]
v0.9 & 2026-06-13 &
  Caché de subsistema en \texttt{KGeometricSIA}: elimina recálculo de tabla BFS
  entre valores de $k$ (ahorro 75\%). \\[4pt]
v0.8 & 2026-06-13 &
  Conteo adaptativo $n_{\mathrm{obj}} = \min(25, \Stir{n}{k})$ y
  warm-start inter-$k$ para sistemas pequeños. \\[4pt]
v0.7 & 2026-05-29 &
  \texttt{letras\_a\_bits}, \texttt{run\_prueba}: utilidades de pruebas manuales. \\[4pt]
v0.6 & 2026-05-29 &
  Análisis experimental: benchmark completo con 4 tablas de datos reales. \\[4pt]
v0.5 & 2026-05-29 & Suite de 13 tests automatizados: grupos A, B, C. \\[4pt]
v0.4 & 2026-05-29 &
  Integración en punto de entrada: variable \texttt{KGEOMIP\_K},
  despacho a \texttt{KGeometricSIA}. \\[4pt]
v0.3 & 2026-05-29 &
  \texttt{KGeometricSIA}: heurística greedy LPT/round-robin/bloques. \\[4pt]
v0.2 & 2026-05-29 &
  \texttt{System.k\_partir()}, \texttt{stirling()}, \texttt{particionar\_conjunto()},
  \texttt{k\_particiones()}, \texttt{fmt\_k\_particion()}. \\[4pt]
v0.1 & --- & Estado inicial: GeoMIP base con soporte exclusivo $k=2$. \\
\bottomrule
\end{longtable}

\section{Referencias y Bibliografía}
\label{apx:referencias}

\begin{enumerate}
  \item Tononi, G. (2004). \textit{An information integration theory of consciousness}.
    BMC Neuroscience, 5(1), 42.

  \item Tononi, G., Boly, M., Massimini, M., \& Koch, C. (2016).
    \textit{Integrated information theory: from consciousness to its physical substrate}.
    Nature Reviews Neuroscience, 17(7), 450--461.

  \item Graham, R.\,L. (1969). \textit{Bounds on multiprocessing timing anomalies}.
    SIAM Journal on Applied Mathematics, 17(2), 416--429.
    (Algoritmo LPT, garantía de aproximación $4/3$.)

  \item Knuth, D.\,E. (1997). \textit{The Art of Computer Programming, Vol.\ 2}.
    Addison-Wesley. (Números de Stirling, Sección 5.2.3.)

  \item Rubner, Y., Tomasi, C., \& Guibas, L.\,J. (2000).
    \textit{The Earth Mover's Distance as a metric for image retrieval}.
    International Journal of Computer Vision, 40(2), 99--121.

  \item Harris, C.\,R., et al.\ (2020). \textit{Array programming with NumPy}.
    Nature, 585, 357--362.

  \item Virtanen, P., et al.\ (2020).
    \textit{SciPy 1.0: Fundamental algorithms for scientific computing in Python}.
    Nature Methods, 17, 261--272.

  \item Cormen, T.\,H., Leiserson, C.\,E., Rivest, R.\,L., \& Stein, C. (2009).
    \textit{Introduction to Algorithms} (3rd ed.). MIT Press.

  \item Mansour, T. (2012). \textit{Combinatorics of Set Partitions}.
    CRC Press. (RGS, números de Stirling.)
\end{enumerate}

% ════════════════════════════════════════════════════════════
%  SECCIÓN KQNodes (PENDIENTE)
% ════════════════════════════════════════════════════════════
\newpage
\section*{Próxima Sección: KQNodes (pendiente)}
\addcontentsline{toc}{section}{KQNodes (pendiente)}

\begin{blueBox}
\textbf{Nota:} La sección correspondiente a \textbf{KQNodes} --- extensión del algoritmo
QNodes (Queyranne) a k-particiones --- se incorporará en la próxima versión de este
documento.

Estructura prevista:
\begin{itemize}
  \item Fundamentos teóricos: submodularidad y algoritmo de Queyranne para $k \geq 3$
  \item Diseño algorítmico de \texttt{KQNodesSIA}
  \item Análisis de complejidad comparado con K-GeoMIP
  \item Resultados experimentales cruzados
  \item Comparativa KGeoMIP vs.\ KQNodes
\end{itemize}
\end{blueBox}

\end{document}
